% Pandoc template for the ACM journal build (JOURNAL_BUILD.py). ACM large format, as DTRAP uses.
% Preprint mode: [acmlarge,nonacm] removes ACM's copyright block and citation footer.
% Anonymous mode: [acmlarge,review,anonymous] hides the author and adds line numbers for review.
\documentclass[acmlarge,review,anonymous]{acmart}
\acmJournal{DTRAP}
\usepackage{longtable,array,calc}
\usepackage{xurl}
\usepackage{seqsplit}
% Reference-list definitions, as Pandoc writes them for the preprint build.
\newlength{\cslhangindent}
\setlength{\cslhangindent}{1.5em}
\newlength{\csllabelwidth}
\setlength{\csllabelwidth}{3em}
\newlength{\cslentryspacingunit} % times entry-spacing
\setlength{\cslentryspacingunit}{\parskip}
\newenvironment{CSLReferences}[2] % #1 hanging-ident, #2 entry spacing
 {\bibliofont % the template's own reference-list size
  % don't indent paragraphs
  \setlength{\parindent}{0pt}
  % turn on hanging indent if param 1 is 1
  \ifodd #1
  \let\oldpar\par
  \def\par{\hangindent=\cslhangindent\oldpar}
  \fi
  % set entry spacing
  \setlength{\parskip}{#2\cslentryspacingunit}
 }%
 {}
\newcommand{\CSLBlock}[1]{#1\hfill\break}
\newcommand{\CSLLeftMargin}[1]{\parbox[t]{\csllabelwidth}{#1}}
\newcommand{\CSLRightInline}[1]{\parbox[t]{\linewidth - \csllabelwidth}{#1}\break}
\newcommand{\CSLIndent}[1]{\hspace{\cslhangindent}#1}
\providecommand{\tightlist}{\setlength{\itemsep}{0pt}\setlength{\parskip}{0pt}}
\setlength{\LTcapwidth}{\textwidth}
\setlength{\emergencystretch}{3em} % as in the preprint build: prevents overfull lines
\begin{document}
\title{Threat Analytics Across Payment Products: Evidence Requirements
for Agent Investigations}
\author{Jonathan Matkowsky}
% No institution. The title line shows "NAME, USA"; the template's automatic commas would print ",, USA",
% and its city and state commands remove the space after the comma, so the address note is written out here.
\affiliation[obeypunctuation=true]{\country{USA}}
\email{Jonathan@matkowsky.com}
\begin{abstract}
Abuse of a payment platform can cross products and accounts, and the
platform's own protective decisions change which outcomes can later be
observed. A blocked payment cannot produce a chargeback; an incomplete
search cannot establish absence; a reused conclusion cannot corroborate
itself. This paper specifies the records an AI agent needs in order to
investigate such cases under human authority. The specification rests on
three requirements. First, each outcome record links to the control
decision that may have shaped it, with separate fields for reporting
delay, outcomes the decision made unobservable, and false reports.
Second, a gateway outside the agent records what each query actually
covered, and every fact retains both the time it happened and the time
the system learned it. Third, a person must release a conclusion before
it enters organizational memory for reuse by later investigations; a
released conclusion serves only as a hypothesis or context, and any
record that an agent output informed must link back to that output.

Draft schemas, 29 cross-record rules and two checkers make the
requirements testable, and adversarial cases written by a second model
probe the rules. Public sources supply the design premises, and Stripe's
public product pages and documentation provide an illustrative product
scope. Two hypothetical investigations show how missing context changes
the conclusions that the same observations can support, replays of two
FTC complaints on synthetic records show what the records expose and
miss, and an evaluation plan proposes tests of evidential quality,
operational cost and agent benefit. The contribution is a joint
specification for investigation across payment products, built on
established methods and recent agent-governance designs. It is a
synthesis, not a new method. The architecture has not been built, no
platform data were used, its only experiments run the record schemas on
synthetic records, and no practitioner has reviewed it.
\end{abstract}
\keywords{payment fraud, abuse investigation, threat intelligence,
selective labels, agent memory, provenance, human oversight}
\maketitle
\hypertarget{introduction}{%
\section{Introduction}\label{introduction}}

An investigation asks whether a card payment attempt drew a fraud
chargeback. An AI agent queries the platform's feed of outcome reports
for the attempt's reporting window, 4 February to 4 June 2026. The query
covers the whole window and returns nothing. The agent drafts the
conclusion that no fraud chargeback arrived, citing the query and a
reading of the chargeback channel that shows its window closed with no
report. A person releases it into the platform's memory, where later
investigations can read it. Table 1 removes one group of records at a
time from this synthetic case.

\begingroup\footnotesize

\begin{longtable}[]{@{}
  >{\raggedright\arraybackslash}p{(\columnwidth - 4\tabcolsep) * \real{0.3000}}
  >{\raggedright\arraybackslash}p{(\columnwidth - 4\tabcolsep) * \real{0.3400}}
  >{\raggedright\arraybackslash}p{(\columnwidth - 4\tabcolsep) * \real{0.3600}}@{}}
\caption{\textbf{One conclusion with each group of records removed.}
Each row is a pair of synthetic histories that the remaining records
cannot tell apart (Supplement S4.4): the first two on this claim, the
third on an earlier conclusion in the same history.}\tabularnewline
\toprule\noalign{}
\begin{minipage}[b]{\linewidth}\raggedright
Records removed
\end{minipage} & \begin{minipage}[b]{\linewidth}\raggedright
What the rest appear to support
\end{minipage} & \begin{minipage}[b]{\linewidth}\raggedright
What the removed records can show
\end{minipage} \\
\midrule\noalign{}
\endfirsthead
\toprule\noalign{}
\begin{minipage}[b]{\linewidth}\raggedright
Records removed
\end{minipage} & \begin{minipage}[b]{\linewidth}\raggedright
What the rest appear to support
\end{minipage} & \begin{minipage}[b]{\linewidth}\raggedright
What the removed records can show
\end{minipage} \\
\midrule\noalign{}
\endhead
\bottomrule\noalign{}
\endlastfoot
The control decision on the attempt and the state of its chargeback
channel (C1) & No chargeback arrived, so the absence counts against
fraud & The platform blocked the attempt, so no chargeback could
arrive \\
The query's coverage (C2) & The query returned nothing, so no chargeback
arrived by the end of the window & The report feed was loaded only to 20
May, so the query did not cover the end of the window \\
The release decision (C3) & The conclusion is in memory, so a later
investigation may use it & No person released it, so it may not be
reused \\
\end{longtable}

\endgroup

Each possibility calls for a different conclusion and a different next
step. The case rests on a short list of records: the control decision on
the attempt, the state of its chargeback channel and a dated reading of
that channel (C1); the query run with its coverage (C2); and the draft,
the release decision and the released item (C3). The broader inventory
of 71 signals and architecture of 24 responsibilities, which connect the
requirements to the evidence an investigator checks and to who checks
it, are reference material in Supplement S3 and S4.

The difficulty is not simply connecting more data. The connections must
preserve what each record can establish. A payment log records activity.
A risk score records a model's assessment. A control decision records
what the platform allowed, blocked or restricted. An outcome label
records a reported result, which may arrive late or be wrong. When these
distinctions disappear, an investigation can mistake a decision for
proof, an unobservable outcome for a benign one, or repeated analysis
for corroboration.

Each of these failures has its own research literature. Selective-labels
and counterfactual-evaluation methods address outcomes hidden by
decisions. Temporal and incomplete-database research addresses when a
fact became known and whether a query answer is complete. Provenance
standards and agent-memory designs address the status of derived
conclusions. These bodies of work mostly treat the failures one at a
time. Some recent agent-memory and decision-record designs address two
or three of them in part, and they differ on who may authorize reuse of
a conclusion and what a reused conclusion can establish (Section 2). An
investigation that spans payment products can meet all three failures in
one case, so its records must keep all three visible together.

This paper specifies records for threat analytics: investigation that
connects product activity to intelligence, proposed detections and
control recommendations, then learns from later outcomes. The setting is
a payment platform offering several products to businesses and their end
customers. The investigator of interest is an AI agent: software that
uses a language model to choose queries and interpret their results. It
works within human-assigned tasks and enforced access limits. It can
investigate and recommend; authorized people or accountable operational
systems decide what happens to accounts, payments and controls.

The design addresses three research questions:

\textbf{RQ1:} What records and analytical capabilities support
investigation across payment products while preserving the meaning of
observations and outcomes?

\textbf{RQ2:} How can an agent test explanations and prepare useful
outputs while people retain review, release and operational authority?

\textbf{RQ3:} How can later outcomes and reusable lessons inform
investigation without concealing the effects of controls or recycling
conclusions as evidence?

The contribution is three requirements on the investigation's records:

\textbf{C1} links outcomes to decisions. An outcome channel is a route,
such as a chargeback, by which an outcome can be reported. Each outcome
record links to the control decision that may have shaped it, and to the
status of each channel that the decision closed, delayed or changed.
Late reports, outcomes that a control made unobservable, and false
reports occupy separate fields. An evaluation can then state the joint
assumptions it needs about delay, observability and report truth, each
of which may depend on the decision (Section 7).

\textbf{C2} records query coverage and learned time. Query coverage is
what a query actually covered: the data and period it examined, and
whether its result was complete. Learned time is when a fact became
available to the system's queries. Every query passes through a gateway,
a checkpoint outside the agent, which records its coverage. Every fact
keeps both when it happened and when the system learned it (Sections
5.1, 5.3 and 6.2).

\textbf{C3} requires human release before promotion into organizational
memory. No agent output enters the evidence plane, which stores
observations, as evidence of the activity under investigation (Sections
4.2 and 5.3); the agent's own outputs are recorded only as observations
of what the agent did (Supplement S4.1). A person releases a conclusion
into organizational memory, where it serves only as a hypothesis or
context. Any record that an agent output informed links back to that
output and counts as dependent, not as independent corroboration; a
record written while an agent output is open to its writer counts as
informed (Sections 4.2, 6.4 and 7.3).

The requirements work together. C1 identifies which outcomes the
platform's decisions left observable. C2 establishes what information an
investigation actually examined and could have known. C3 prevents the
resulting interpretation from returning as independent support for
itself. A record can satisfy one requirement and still mislead an
investigation through a failure of another (Section 8.3).

Draft schemas, 29 cross-record rules and two checkers make the
requirements testable, and every rule has a failing case (Supplement
S4). Fifteen histories written by a second model to test the rules
exposed eight gaps; after revision on those same histories, the rules
reject nine of the 13 meant to break a requirement, and the other four
are stated limits. Two hypothetical investigations show which
conclusions survive when particular records are missing, and replays of
two FTC complaints on synthetic records show what the records expose and
miss (Section 8).

The intended readers design, audit or evaluate investigation systems.
Stripe's public product pages and documentation supply an illustrative
product scope, not a description of its internal systems. The paper
synthesizes public evidence and proposes a design. It reports no
implementation of the architecture and no measurement on platform data;
its only experiments run the record schemas on synthetic records
(Sections 8.4 and 8.5). Section 2 establishes its relationship to prior
work and Section 3 explains how it was derived. Sections 4 and 5 answer
RQ1, with Section 6.3 and Supplement Table S4.2 adding the agent's
analytical capabilities; Section 6 answers RQ2 and Section 7 answers
RQ3. Section 8 shows the consequences of the design in two hypothetical
investigations, and Section 9 states how its value would be tested.

\hypertarget{relationship-to-prior-work}{%
\section{Relationship to prior work}\label{relationship-to-prior-work}}

The three requirements combine established ideas in a specific setting.
This section identifies the closest precedents for each requirement and
the choices the specification adds. Supplement S9 gives the extended
comparison, with the nearest precedents for each requirement summarized
in Table S9.1; S1 provides source passages, publication status and a
50-source requirements matrix.

\hypertarget{decisions-shape-the-outcomes-used-to-judge-them}{%
\subsection{Decisions shape the outcomes used to judge
them}\label{decisions-shape-the-outcomes-used-to-judge-them}}

When a decision determines which cases receive outcome labels,
evaluation faces the selective labels problem
{[}\protect\hyperlink{ref-Lakkaraju2017Selective}{1},
\protect\hyperlink{ref-Kilbertus2020Fair}{2}{]}. Counterfactual analysis
of logged decisions is most direct when the decisions were randomized
with known probabilities. Each logged sample is then reweighted to
reflect its probability under another decision rule, so that average
outcomes under that rule can be estimated
{[}\protect\hyperlink{ref-Bottou2013Counterfactual}{3}{]}. All of these
methods require assumptions about which comparisons the data support
{[}\protect\hyperlink{ref-HernanRobins2020}{4}{]}.

Payment research makes the reporting problem concrete. TitAnt trains on
delayed user fraud reports
{[}\protect\hyperlink{ref-Cao2019TitAnt}{5}{]}. Dal Pozzolo and
colleagues distinguish prompt investigator feedback from later labels
and show how alert-driven review biases training data
{[}\protect\hyperlink{ref-DalPozzolo2017Fraud}{6}{]}. An e-commerce
research agenda likewise separates suspicion on blocked transactions
from outcome labels {[}\protect\hyperlink{ref-Tax2021Fraud}{7}{]}. A
theoretical payments preprint models authorization, whether fraud is
reported, whether the report has arrived by training time, and label
corruption, for one reporting channel
{[}\protect\hyperlink{ref-Dhama2026Causal}{8}{]}. Related preprints
examine delayed, incomplete or incorrect feedback
{[}\protect\hyperlink{ref-Dhama2026Limits}{9},
\protect\hyperlink{ref-HeLoh2026Grab}{10}{]}. Chapelle's advertising
model supplies a precedent for treating a click with no conversion yet
as unlabeled rather than negative
{[}\protect\hyperlink{ref-Chapelle2014Delayed}{11}{]}. Two recent design
preprints come closest to C1's record. One, the Decision Event Schema,
seals each decision record once its hash is computed. Facts learned
later, such as the arrival of ground truth, are appended as linked
records that reference the decision, and links between decisions carry a
type {[}\protect\hyperlink{ref-Solozobov2026DecisionTrace}{12}{]}. The
other, an execution-coverage model, links records to decisions and
distinguishes refused execution branches from late reports
{[}\protect\hyperlink{ref-Wu2026EvidenceCoverage}{13}{]}.

C1 applies these ideas to several outcome channels at once. It records
what each decision closed, delayed or changed, separately for each
channel and read time. A blocked card payment cannot produce a
chargeback on that attempt, but an incident finding remains possible. C1
keeps the decision as the unit of record, as the Decision Event Schema
does. Each outcome report is a separate appended record. It links to
every decision that may have shaped it, and each link states that
decision's role. The schema has no field for outcome channels, but C1's
fields could extend it rather than replace it (Section 7.1; Supplement
S4.2). The proposed records expose the assumptions needed to interpret
such outcomes; they do not supply a validated estimator. Section 7
develops this distinction.

\hypertarget{evidence-depends-on-coverage-time-and-origin}{%
\subsection{Evidence depends on coverage, time and
origin}\label{evidence-depends-on-coverage-time-and-origin}}

Temporal databases distinguish when a fact held from when it was stored
{[}\protect\hyperlink{ref-Snodgrass1986Temporal}{14}{]}, and
agent-memory designs retain the same two kinds of time
{[}\protect\hyperlink{ref-Wang2026TOKI}{15}--\protect\hyperlink{ref-Dan2026Anamnesis}{17}{]}.
The execution-coverage model of Section 2.1 also separates event time
from verifier receipt time, and it judges whether a verifier's evidence
was complete at a stated cutoff from the records then due
{[}\protect\hyperlink{ref-Wu2026EvidenceCoverage}{13}{]}. Research on
incomplete databases relates the completeness of a query answer to which
data are known to be complete
{[}\protect\hyperlink{ref-Razniewski2011Completeness}{18}{]}. These
precedents matter whenever a query returns no matching records or a
relationship becomes known after its effective date.

C2 combines explicit query coverage with learned time. The
execution-coverage model asks whether the records an execution owes had
arrived by a cutoff. C2 records what each investigative query examined.
A gateway outside the agent records the requested and examined data,
access limits and known gaps. Each fact's learned time is when it became
available to queries. A relationship learned today cannot appear in
yesterday's reconstructed decision merely because it was in effect then,
although later outcomes may still evaluate that fixed decision under a
stated outcome window.

The W3C PROV data model expresses the derivation of one entity from
another {[}\protect\hyperlink{ref-W3C2013PROVDM}{19}{]}, and decision
provenance records the data flows surrounding decisions
{[}\protect\hyperlink{ref-Singh2019DecisionProvenance}{20}{]}. STIX
separates observed data from intelligence assertions and opinions
{[}\protect\hyperlink{ref-OASIS2021STIX}{21}{]}, and US intelligence
analytic standards ask analytic products to distinguish underlying
information from assumptions and judgments
{[}\protect\hyperlink{ref-ODNI2015ICD203}{22}{]}. For agent
investigations, those distinctions must survive retrieval and reuse. The
agent's own account of its work cannot establish what its tools actually
examined; independently written run records are the relevant precedent
{[}\protect\hyperlink{ref-Dai2026Hearsay}{23}{]}.

\hypertarget{released-conclusions-remain-conclusions}{%
\subsection{Released conclusions remain
conclusions}\label{released-conclusions-remain-conclusions}}

Truth-maintenance research records the support for conclusions and
revises them when that support changes
{[}\protect\hyperlink{ref-Doyle1979TMS}{24},
\protect\hyperlink{ref-Malheiro1994BeliefRevision}{25}{]}. A conceptual
framework for language agents separates working, episodic, semantic and
procedural memory {[}\protect\hyperlink{ref-Sumers2024CoALA}{26}{]}, and
a survey argues that agent memory should keep raw observations,
extracted facts and inferred items apart, with derivation links
{[}\protect\hyperlink{ref-Wang2026AgentTraces}{27}{]}. Together these
establish substantial precedent for C3.

The closest designs differ in who authorizes reuse and what a reused
item can establish. Anamnesis starts model-derived claims as candidates,
and no memory operation can raise their authority without a separately
recorded grant; the grantor need not be a human
{[}\protect\hyperlink{ref-Dan2026Anamnesis}{17}{]}. GovMem groups
correlated agent traces by their dependencies when assessing support and
permits automatic promotion or human review
{[}\protect\hyperlink{ref-Qi2026GovMem}{28}{]}. Quipu's write gate
requires human approval only under some policies
{[}\protect\hyperlink{ref-Brown2026Quipu}{29}{]}. A data-governance
preprint is closer on authorization: no agent-proposed metadata decision
is committed before a steward reviews it, and a proposed gate admits
only verified records as precedent, records each promotion and lets a
reviewer demote a record, flagging the decisions that relied on it. It
does not say who authorizes promotion or what a retrieved precedent may
establish {[}\protect\hyperlink{ref-Akhtar2026Contestable}{30}{]}. In
Stripe's agent-assisted fraud modeling, engineers review the agents'
hypotheses and results and decide what to deploy, and for one fraud
attack the agents were given results from earlier experiments and
investigations; the account does not describe a release step for those
earlier results
{[}\protect\hyperlink{ref-Drapeau2026ScalingRadar}{31}{]}. A
threat-detection design comes closest on reuse as evidence. It tags each
detection compiled from one of its hypotheses and drops tagged
observations at ingest, so the system's own conclusions cannot re-enter
as fresh evidence. It leaves as a residual risk an analyst who writes
such an alert into a report that is later ingested, and its reference
implementation does not yet exercise the rule
{[}\protect\hyperlink{ref-Evron2026AdversaryBrains}{32}{]}.

C3 adopts human release for every conclusion entering organizational
memory. A released conclusion is reused only as context or a hypothesis,
a limit the data-governance design does not state for its precedents;
release does not turn it into an observation. Records informed by an
agent output retain their dependence on it, including labels, detections
and control decisions. Because the interfaces that write these records
add the links, they also cover a person who writes a record while an
agent output is open in the case, part of the relay the threat-detection
design leaves open; a missing link still does not establish independence
(Section 4.2). This is a design choice whose benefit and cost require
testing, particularly through the agent-benefit and release-gate
evaluations (E5 and E8, Section 9.1); it is not a claim that human
approval guarantees truth.

\hypertarget{what-the-combined-specification-contributes}{%
\subsection{What the combined specification
contributes}\label{what-the-combined-specification-contributes}}

Industrial payment systems supply examples of fast decisions and
relationships between accounts
{[}\protect\hyperlink{ref-Cao2019TitAnt}{5},
\protect\hyperlink{ref-Liu2018GEM}{33}{]}. Payment providers' public
documentation already exposes parts of the first two requirements.
Stripe lets businesses query which rules affected each charge and join
disputes to charges
{[}\protect\hyperlink{ref-Stripe2026QueryDisputes}{34}{]}. Its dispute
analytics warns that rates by payment date can look lower for recent
payments because customers can wait weeks to dispute
{[}\protect\hyperlink{ref-Stripe2026DisputeAnalytics}{35}{]}, and each
scheduled query run records the as-of time of the data snapshot it
queried {[}\protect\hyperlink{ref-Stripe2026SigmaRun}{36}{]}. Another
provider's risk dashboard shows whether transactions its risk engine
allowed later resulted in a chargeback, and lists separately those it
blocked {[}\protect\hyperlink{ref-Adyen2026RiskMonitoring}{37}{]}. The
specification states what such joins, labels and timestamps must
preserve for an investigation: which outcomes each decision closed,
delayed or changed, and which reports are late or false. A governance
synthesis for fraud risk decision systems combines decision traces,
measures of evidence sufficiency under delayed labels and monitoring
without labels, and extends them analytically to multi-agent systems
{[}\protect\hyperlink{ref-Solozobov2026Synthesis}{38}{]}. It rates
architectures by how many decision-record properties they fill by
construction; it specifies no outcome channels, query coverage record or
release rule for reuse. Supplement Table S1.1 locates these capabilities
across research families.

Financial-sector and audit documents on AI, including agents, address
adjacent questions. They include a handbook from Project MindForge, an
industry initiative led by the Monetary Authority of Singapore
{[}\protect\hyperlink{ref-MindForge2026Operationalisation}{39}{]}, a
white paper on runtime safeguards for agentic finance
{[}\protect\hyperlink{ref-MAS2026SAFR}{40}{]} and guidelines from
India's computer emergency response team for auditing AI systems
{[}\protect\hyperlink{ref-CERTIn2026AIAudit}{41}{]}. Their provisions on
logging tool calls, on AI analysis as a working aid and on measuring
human review partly parallel C2, C3 and E8 (Supplement S9.3). None of
the three documents asks for outcome records by channel (C1), separates
when a fact held from when it became available to queries (C2), or
requires human release before a conclusion is reused from memory (C3).

Two runtime-governance designs for financial agents bar context from
authorizing actions and limit what counts as evidence; neither requires
human release before reuse
{[}\protect\hyperlink{ref-Tang2026SAGEFin}{42},
\protect\hyperlink{ref-Szpruch2026RuntimeGovernance}{43}{]}.

Two payment-agent preprints show why specification and operational
benefit must be kept apart. SR-Fraud separates scoring from reflection
and tests proposed rules before reuse. On one replay benchmark it
reports better detection than its scoring agent alone
{[}\protect\hyperlink{ref-Tan2026SRFraud}{44}{]}; Supplement S9.2 gives
its comparisons with classifiers. On 60 balanced synthetic cases that
its classifier scored between 0.2 and 0.8, Sharma's agent was correct in
65.0 percent, below the 71.7 percent from thresholding the same
classifier score at 0.5. Six of its eight overrides turned a correct
decision into an error
{[}\protect\hyperlink{ref-Sharma2026AuditableFraud}{45}{]}. Neither
result establishes general agent performance: one rests on a single
non-public replay benchmark, and the other on a small, balanced
synthetic sample. Together they motivate comparing an agent with simpler
alternatives given the same information.

This paper connects decision-shaped outcomes, query coverage and
controlled reuse in one payment-investigation specification. The traces
show which inferences lose support when those records are removed; they
establish logical consequences within hypothetical cases, not improved
detection. The focused comparison does not establish uniqueness or
exhaustive coverage, and operational value depends on the evaluations in
Section 9.

\hypertarget{deriving-the-specification-from-public-evidence}{%
\section{Deriving the specification from public
evidence}\label{deriving-the-specification-from-public-evidence}}

The design starts from documented product capabilities and asks how an
adversary could misuse each one. For each abuse mechanism, meaning a way
abuse can occur, it then identifies the records needed to distinguish
the malicious explanation from plausible alternatives. Stripe's public
product catalog, roadmap and documentation define the illustrative
scope; research, engineering accounts, incident reports and guidance
supply the design premises.

The scope, read on 27 September 2026, contains 88 product and feature
entries. A retained mechanism must target a documented capability, serve
a plausible adversary goal and have supporting public evidence. Public
job postings from Stripe for abuse, risk and machine-learning roles
helped prioritize topics as a statement of its current priorities, not
as evidence of incidents {[}\protect\hyperlink{ref-Stripe2026}{46},
\protect\hyperlink{ref-StripeARE2026}{47}{]}. Supplement S2 records the
99 candidate mechanisms, how postings set priorities, the nine search
paths and each candidate's final disposition; every mapped mechanism
also has supporting sources beyond postings.

The resulting inventory, 71 signals in 12 families linked to 61 of the
99 candidate mechanisms, is in Supplement S3. Supplement S2 gives the
disposition of the other 38 candidates.

Three published schemas describe the platform's application programming
interface (API), the interface through which software requests
operations and receives data. Supplement S5 supplies the object, field
or event supporting each match.

Supplement S1 records the version, relevant passages and use of each of
the 127 works cited in the manuscript and supplement, and compares 50
sources against the three requirements. Supplement S3 lists the 61
design sources and the depth of each read. These records let readers
inspect the synthesis; their counts do not measure evidence quality.

The synthesis distinguishes sourced findings, design choices and
hypotheses (Supplement S2.1). One author, assisted by AI models, made
the mappings and coverage judgments; agreement between independent human
coders has not been measured. A model that took no part in the mapping,
given a linking rule written afterward, made 49 of 52 sampled links and
70 that the paper lacks. The links are therefore a sourced subset, not a
complete map (Supplement S3). Section 9.2 explains the limits this
focused selection places on coverage and transfer.

The author used Anthropic Claude and OpenAI ChatGPT/Codex in the
research. The models searched for and read the sources, proposed the
mappings and comparison codes, ran the prior-work search, wrote the
scripts, synthetic data and figures, and reviewed drafts. Every cited
work was read in full, or, for a book or standard, the cited chapters;
Supplement S1 records the version and passages for each and names the
few parts not read, such as reference pages. The author read each cited
work to the same depth, and the manuscript and supplement in full, more
than once before this version was released. A check script recomputes
the counts, lineage and arithmetic on every proposed change. The author
decided what entered the paper. Model review is not independent human
peer review.

\hypertarget{from-product-activity-to-protective-decisions}{%
\section{From product activity to protective
decisions}\label{from-product-activity-to-protective-decisions}}

This section identifies the parties whose activity the records describe,
the path from records to protective decisions along which meaning must
be preserved, and the responsibilities that path requires. It begins the
answer to RQ1, which Section 5 continues and Section 6.3 completes.

\hypertarget{the-platform-and-its-participants}{%
\subsection{The platform and its
participants}\label{the-platform-and-its-participants}}

The proposed payment platform serves businesses, which serve end
customers. A business holds a platform account. It may instead hold a
connected account under another business that has integrated the
platform into its product. A merchant is a business accepting payments
through either account type. People act for businesses through operator
accounts. Software agents may use delegated credentials for a business
or end customer; these are distinct from the investigating agent.

These roles determine whose activity a record describes. An end
customer's action is not automatically an action by the business or its
operator.

\hypertarget{the-path-from-records-to-action}{%
\subsection{The path from records to
action}\label{the-path-from-records-to-action}}

Data pipelines collect telemetry and state records, map them to shared
entity and field definitions, and preserve their origins. Storage and
correlation services make them searchable by time and entity.
Investigators use them to prepare intelligence, meaning assessments
intended to inform decisions, together with proposed detections and
recommendations. Authorized people or accountable operational systems
decide whether to change controls, and those changes affect what later
records can show (Figure 1).

Three feedback loops connect this work. \textbf{Investigative feedback}
changes the next hypothesis or test within a case.
\textbf{Organizational-learning feedback} makes released intelligence or
tested procedures available to later cases as context.
\textbf{Protective feedback} changes the platform through an authorized
control or product decision. Their outputs require different authority
and support different inferences.

The architecture separates three logical layers, called planes. They
need not use separate databases:

\begin{itemize}
\tightlist
\item
  The \textbf{evidence plane} holds observations with their source, time
  and entities. Scores and control decisions belong here as observations
  of what a model concluded or a control did. Derived features belong
  here too, as values recomputable from recorded inputs; a feature
  computed over a cluster cites the cluster version it used, and the
  cluster itself remains a threat-knowledge item. A feature computed
  over a cluster that an agent proposed, or that depends on an agent
  output or released item, links back to that output or item and counts
  as dependent on it.
\item
  The \textbf{structured threat-knowledge plane} holds revisable
  inferences, such as relationships between entities, groups of activity
  and conclusions with confidence and provenance.
\item
  The \textbf{organizational-memory plane} holds released intelligence
  and tested lessons about procedure. These guide later work as
  hypotheses or context.
\end{itemize}

\begin{figure}
\centering
\includegraphics[width=\linewidth,height=0.5\textheight,keepaspectratio]{FIGURES/FIG1_LOOP.pdf}
\Description{Figure 1: Product records, investigation, review, authorized decisions and feedback}
\caption{\textbf{From product records to protective decisions.} Investigation
produces drafts. In the normal path, shown here, independent quality
review precedes human release. Under Section 6.4, an authorized human
may also decide urgent containment before review, without releasing the
draft into organizational memory. Each resulting control decision
records what the platform did, not whether the suspected abuse occurred.}
\end{figure}

For example, two accounts using the same device is an observation.
Common control is an inference. A released report retaining that
qualified inference becomes organizational memory; retrieving it creates
no new observation. An investigator may follow its lineage to cite the
original observations at their original times and sources. A
threat-intelligence reference model likewise traces the creator and
modifiers of each intelligence product
{[}\protect\hyperlink{ref-Sakellariou2022CTIReferenceModel}{48}{]}.

Dependence must also survive indirect reuse. A human label, assessment,
control decision, detection or threat-knowledge item informed by an
agent output links back to that output. Labeling, assessment, control,
detection and threat-knowledge interfaces automatically attach every
agent output and every released organizational-memory item open in the
writer's case or session. Only a person with release authority may
remove a link, and removal is logged. Intelligence sent outside the
platform carries an identifier. When content bearing that identifier
returns through an external feed, the external-knowledge record links to
the released item and counts as dependent on it.

These records remain dependent on the agent output that informed them.
They cannot independently corroborate the agent's findings, a detection
derived from its work, or the same hypothesis. The links capture only
exposure observed by the interfaces: an unlinked record may still
reflect something its writer read elsewhere. A missing link therefore
does not establish independence. Detections that rely on agent-dependent
labels are validated on external outcome channels or on cases worked by
staff without access to agent outputs.

This separation cannot repair missing source data. An agent cannot
retrieve telemetry never collected, reconstruct timing discarded by a
pipeline, or establish an entity link without suitable identifiers.

\hypertarget{responsibilities-and-authority}{%
\subsection{Responsibilities and
authority}\label{responsibilities-and-authority}}

Supplement S4 defines 24 responsibilities for the investigation system.
They are functions, not proposed teams or a required number of agents.

Payment authorization, payouts and account-review systems produce C1's
decision records; their operational owners sit outside the investigation
responsibilities.

Human tasking and release stay with people, and quality review remains
independent of the work checked. Other functions may use one agent with
tools, separate model sessions, several agents, fixed rules or a
combination. Additional separation must justify its cost through
narrower authority, access or failure scope.

The investigating agent may query permitted records and prepare
intelligence, hunting and detection candidates, recommendations and
proposed lessons. It cannot change accounts, payments or controls,
certify its own output, or release conclusions into organizational
memory.

\hypertarget{records-an-investigator-can-reason-from}{%
\section{Records an investigator can reason
from}\label{records-an-investigator-can-reason-from}}

This section answers the records part of RQ1 in four steps: record
classes that state what each record can establish (5.1), the abuse
mechanisms that motivate particular records (5.2), the properties each
store must preserve (5.3), and what published schemas already supply and
privacy obligations permit (5.4 and 5.5). Section 6.3 adds the
capabilities.

\hypertarget{separate-observations-from-conclusions}{%
\subsection{Separate observations from
conclusions}\label{separate-observations-from-conclusions}}

A payment log, risk score and fraud report may concern the same
transaction but establish different facts. The log records activity, the
score a model's assessment, and the report what someone alleged. Copies
of one report do not provide several independent confirmations. The
record's class makes these limits explicit.

The same principle applies to relationships. A shared payout destination
is observed; common control is inferred. Attribution, identifying a
particular actor behind the activity, is a stronger claim requiring
further support. An earlier attribution remains historical context when
retrieved in a new case.

The specification uses seven record classes (Supplement Figure S3.1):

\begin{itemize}
\tightlist
\item
  \textbf{Source observations} record events or reports at their source.
\item
  \textbf{Normalized observations} map source records to common entity
  and field definitions while retaining provenance.
\item
  \textbf{Derived features} compute values from recorded inputs for
  specified entities and time windows.
\item
  \textbf{Model or risk outputs} record assessments and the identity and
  version of the model.
\item
  \textbf{Control outcomes} record decisions, independently of whether
  those decisions were correct.
\item
  \textbf{Analytical inferences} state conclusions drawn from evidence.
\item
  \textbf{Finished intelligence}, also called released intelligence, is
  analysis approved for reuse.
\end{itemize}

External knowledge is identified separately by source and type. It may
be an observation, such as an exposed API key, an assessment, or a
published catalog of abuse techniques. It supports what that source
observed or described, not by itself a claim that abuse occurred inside
the platform.

A class describes meaning, not reliability. A derived feature is only as
reliable as its inputs and computation. An outcome label records who
reported what and when; any separately established finding remains
distinct.

Each normalized observation preserves its source and event identifiers,
schema version, entities and lineage. It records event time and learned
time, plus actor and request identifiers where available. Its
environment distinguishes production from tests, synthetic data, replay
or emulation. Its \emph{tenant scope} identifies the business or other
boundary within which the data may be accessed and used. Delivery
identifiers prevent repeated delivery from counting or acting on one
event twice.

Supplement Table S3.1 organizes the signals by family and class.
Supplement S3 defines each signal and the inputs of every derived
feature.

\hypertarget{how-abuse-mechanisms-motivate-the-records}{%
\subsection{How abuse mechanisms motivate the
records}\label{how-abuse-mechanisms-motivate-the-records}}

Each group of abuse mechanisms motivates particular records; Supplement
S3.1 gives the public evidence for each and its limits. This section
keeps the records for outcome labels, on which C1 depends.

\hypertarget{labels-need-their-reporting-path}{%
\subsubsection*{Labels need their reporting
path}\label{labels-need-their-reporting-path}}
\addcontentsline{toc}{subsubsection}{Labels need their reporting path}

An outcome label can be missing because its report is late or because a
control prevented the outcome. It can also be present and false.
Payment-specific theoretical work distinguishes these problems and
proposes estimation methods under assumptions
{[}\protect\hyperlink{ref-Dhama2026Causal}{8}{]}. The specification uses
those distinctions to define records, without assuming that an estimator
recovers hidden platform outcomes.

Each report retains its type, source, lag and attributed transaction or
entity. It keeps three times: when the reported outcome occurred, where
known; when the report was made at its source; and when it arrived at
the platform. Its channel supplies a reporting window and clock-start
event. The evaluation separately records its chosen maturity period, how
long it waits before reading a missing report as a negative. A period
shorter than the channel's window, or any period for a channel without
one, rests on a delay model (Section 7.1). Every decision that may have
shaped the report remains linked, and there can be several. A dispute
can follow a blocked attempt, a retry and an approval, and one account
restriction can affect many payments. Each link therefore states the
decision's role. The report's later status, such as withdrawn, reversed,
won or lost, has its own history, which records whether each won or lost
status followed adjudication. A status change or correction is a new
record that references the earlier one, so what the platform believed at
each time can be reconstructed. Some channels never adjudicate a report,
yet a provider may still record it as lost: Stripe's dispute guide says
disputes that network rules or local regulations make unchallengeable
are closed as lost when they open
{[}\protect\hyperlink{ref-Stripe2026Disputes}{49}{]} (Section 7.1).

An estimate that a report is false belongs to the assessment store,
linked to the report with the responsible model or reviewer and version.
This matters when supporting material itself may be fabricated, as
participants describe for refund fraud involving generative AI
{[}\protect\hyperlink{ref-Zhang2026Refund}{50}{]}. A report may be both
late and false, so one status cannot replace the other.

\hypertarget{preserve-meaning-across-stores-time-and-relationships}{%
\subsection{Preserve meaning across stores, time and
relationships}\label{preserve-meaning-across-stores-time-and-relationships}}

The investigation must be able to retrieve an observation, its
interpretation and the decision it informed separately. These
distinctions concern logical records; they do not require separate
database products. Supplement Table S4.1 lists the properties each store
must preserve.

\textbf{Write authority follows the record's meaning.} Ingestion appends
to the evidence plane. Correlation software and the agent may propose
threat-knowledge items; a person or accountable process accepts them.
Agent proposals remain marked context and cannot support current-conduct
claims without their underlying observations. Only a person releasing a
reviewed draft or a tested, reviewed lesson adds or changes a claim in
organizational memory. Later status records for a released item, such as
monitoring findings, expiry and withdrawal, add no claim and record who
or what wrote them. Rule checkers and quality reviewers write their
findings on each draft version to the agent-and-approval ledger.

\textbf{Event time and learned time answer different questions.} A
relationship added today with an earlier effective date was unavailable
to yesterday's decision. Revisions, external facts, graph membership and
feature definitions therefore retain the versions used. A derived fact's
learned time is when that result became available to queries, not merely
the latest availability of its inputs. It cannot precede any input's
learned time.

A decision's reproducible snapshot need not copy the entire store. The
system logs the inputs used and retains versions from which other
then-available inputs can be reconstructed. This supports checking both
what the decider used and what could have been known.

\textbf{Evaluation time differs from decision time.} Later outcomes may
evaluate an earlier prediction whose features remain fixed. The
evaluation declares its outcome window and maturity period. Training and
model selection use a separate cutoff: labels unavailable then cannot
enter them. Applying the decision-time cutoff to evaluation labels would
wrongly exclude the late outcomes being measured.

\textbf{Channel state differs from report presence and access.} For each
affected item, channel and read time, the records distinguish an open
channel, a temporary delay, a permanent closure and unknown state. A
missing report matures only after a defined window that remained open
and was fully covered. If a decision changed the channel's meaning, as a
liability shift after successful authentication does (Section 7.1), a
matured absence is read under that changed meaning and does not by
itself support a negative for the original outcome. A channel without a
defined window never supports a negative by maturity alone. A matured
absence is provisional, not final. Network and bank rules can reopen or
extend a window after it has elapsed, as when a monitoring program lets
issuers submit fraud disputes past the standard deadline
{[}\protect\hyperlink{ref-Stripe2026MonitoringPrograms}{51}{]} or banks
accept a late return
{[}\protect\hyperlink{ref-Stripe2026ACHTerms}{52}{]}. Each such rule
event is recorded as a channel-state transition with its source, rule
version and learned time, although no platform decision caused it. A
decision's effect on each channel is fixed when the decision is
recorded; later transitions are separate records that never edit it.
Inaccessible records remain unavailable without erasing a known channel
state or an earlier report. Supplement S4.1 specifies the clock-start
and transition fields, with further rules on program status, revised
negatives and reopened reports.

\textbf{An entity link is not a claim of common control.} A verified
identifier from the same namespace and relevant interval supports a
stronger, or hard, link. A weaker, or soft, link may have several
explanations: benign accounts can share a proxy IP address
{[}\protect\hyperlink{ref-Liu2018GEM}{33}{]}. Both retain source and
confounders, meaning plausible alternative explanations. Account-age
records address aged accounts sold for abuse
{[}\protect\hyperlink{ref-Thomas2013}{53}{]}. Where crypto products are
in scope, wallet-to-spender links record permissions motivated by
phishing that induces harmful transactions
{[}\protect\hyperlink{ref-Chen2025Phishing}{54}{]}. Supplement S4
defines each relationship's fields.

\hypertarget{what-public-schemas-reveal}{%
\subsection{What public schemas
reveal}\label{what-public-schemas-reveal}}

This comparison shows how much of a required record a published
interface already supplies and where the specification goes beyond it.
It asks whether a proposed observation has at least a partial
counterpart in a published object, field or event. Where the question
applies, most signals have a partial counterpart, and a keyword search
of the three published schemas found none for the rest.

Supplement S3 records the classification and its counts; S5 provides
schema pins, hashes, event extracts and exact field locations. These
findings describe the inspected schemas, not Stripe's private telemetry
or every possible public interface.

\hypertarget{privacy-and-operating-constraints}{%
\subsection{Privacy and operating
constraints}\label{privacy-and-operating-constraints}}

Privacy and operating obligations constrain the specification;
Supplement S4.1 states them in full. A deployment must specify retention
and deletion for each record class before collecting it. Append-only
logs must still permit required deletion. The agent-and-approval ledger
uses references to regulated data rather than copies. Cross-business
linking requires a stated purpose and authorization for each identifier
type. This paper does not determine which legal basis applies in a
particular deployment.

\hypertarget{how-the-agent-investigates}{%
\section{How the agent investigates}\label{how-the-agent-investigates}}

This section answers RQ2. It follows one investigation from its
assignment to the release of its outputs, then addresses the inputs that
could subvert it. Section 6.3 and Supplement Table S4.2 also complete
the answer to RQ1 by listing the agent's analytical capabilities with
the limit on what each can establish.

\hypertarget{set-the-question-and-scope}{%
\subsection{Set the question and
scope}\label{set-the-question-and-scope}}

A human assigns an investigative question or approves a standing task
template. The assignment defines permitted data, entities, period,
current period and outputs. Each run records the assignment, or the
template and its version, that authorized it, and duplicate alerts are
merged into cases before tasking.

The agent applies the record-class limits while reasoning. A score can
direct attention without establishing abuse. A control outcome records a
decision of potentially unknown correctness. A feature retains its
inputs, computation period and data freshness. A label must be read with
its reporting path and the decisions that shaped it.

Earlier intelligence may suggest a hypothesis or a useful first query.
Retrieval returns its publication date, relevance period, underlying
sources, confidence, restrictions and corrections where recorded.
Several reports repeating one observation remain one source of evidence.
Claims about current conduct require observations from the current
period.

\hypertarget{test-explanations-and-interpret-the-query-that-ran}{%
\subsection{Test explanations and interpret the query that
ran}\label{test-explanations-and-interpret-the-query-that-ran}}

An investigation asks which observation would distinguish plausible
explanations. The agent converts that question into a permitted query,
such as reconstructing a timeline or following an entity link. It first
establishes what the available schemas and services mean.

The loop is \textbf{context → hypothesis → permitted test → observation
→ interpretation → revised hypothesis or conclusion}. Its stopping point
may be an inconclusive assessment when the observations cannot separate
the remaining explanations.

A complete result with no rows establishes only that no accessible rows
matched. It counts against a hypothesis when the suspected activity
should have produced that observation and the records cover where it
would appear. Other results require different interpretations:

\begin{itemize}
\tightlist
\item
  A partial result supports no conclusion about the uncovered data or
  period.
\item
  A truncated result may contain useful matches but cannot establish a
  population total.
\item
  A denied query records an authorization decision.
\item
  An unavailable dataset means the intended test did not occur.
\end{itemize}

The gateway that enforces the agent's scope (Section 6.5) also attaches
the coverage record, independently of the agent. It records the
requested and evaluated entities, with their tenant and identifier
namespaces, the requested and evaluated event intervals, source and
schema versions and any sampling. It also records the as-of cutoff,
which limits the query to facts the system had learned by that time. For
each source it records the event and ingestion watermarks, meaning how
far event time and loading have progressed, known gaps and the retention
horizon. A requested interval that extends past a source's event
watermark, begins before its retention horizon or overlaps a known gap
leaves that source's result partial, as does sampling of that source. It
also records whether access filtering removed rows or row, time or cost
limits were reached.

Fixed rules derive complete, partial, truncated, denied or unavailable
status from that record; several statuses may apply. The agent can
interpret results but cannot alter coverage fields. An unrecorded gap
can still make a result appear complete, so recorded completeness is not
knowledge of every possible loss.

The rule checker (Section 6.4) and the quality reviewer read the full
coverage record. The agent sees only information allowed by its scope,
including counts only when disclosure is authorized. A scope-limited
flag identifies when authorized scope is narrower than the task's named
entities, without disclosing outside records. The checker restricts
absence claims accordingly. An unassessable hypothesis remains distinct
from a query's status.

\hypertarget{capabilities-and-their-boundaries}{%
\subsection{Capabilities and their
boundaries}\label{capabilities-and-their-boundaries}}

Supplement Table S4.2 lists the agent's permitted capabilities. Every
use is scoped, limited and logged; none grants operational authority.

Entity links and labels require particular care. The agent retrieves
each cluster's membership rule and version and tests benign
explanations, including shared proxies
{[}\protect\hyperlink{ref-Liu2018GEM}{33}{]}.

Likewise, no chargebacks may mean no abusive completed payments, blocked
attempts, immature reporting or reports attributed elsewhere. Resolving
a pre-dispute alert, an issuer notice that can be addressed before a
chargeback, may also avert one. The agent examines decision gates,
report times and transaction links before interpreting the absence.

A detection candidate states the behavior to detect, its inputs and an
evaluation plan. That plan covers label provenance and how an adversary
might adapt.

\hypertarget{check-the-claims-then-authorize-their-use}{%
\subsection{Check the claims, then authorize their
use}\label{check-the-claims-then-authorize-their-use}}

A draft separates material claims from context. A claim is material when
its falsity would change a recommendation, detection candidate or
released statement about named entities. Each material claim cites
supporting records with their classes, event intervals and root
observations, reached through lineage. It also lists the retrieved
records that count against it, and quality review adds any it finds. The
checker also identifies every record that the case's query runs returned
within the claim's scope but that the draft neither cites nor lists
against it; Supplement S4 gives how these records are shown and sampled.
A claim that something is absent also cites the query runs whose
coverage records support it. A claim that states or relies on an outcome
that was not reported also cites, for each affected item, the state of
each relevant outcome channel at the read time. The absence counts as a
negative only where the channel matured with no report under its
original meaning, and that negative remains provisional (Section 5.3).
Claims about current conduct cite at least one source observation,
normalized observation or derived feature whose event interval falls
entirely within the task's current period. A score or control decision
from that period does not meet this rule by itself, because it records
an assessment or an action, not the conduct. A derived feature that
carries an agent-output or released-item link does not meet the rule by
itself either; the claim must also cite, for each named entity, a
current-period observation among that feature's inputs.

A checker outside the agent enforces these rules before quality review.
It rejects substitution of organizational memory for observations. It
counts repeated records of one root observation once. That count is a
lower bound: several extractions from one root can add information, but
no more than the root holds
{[}\protect\hyperlink{ref-Bara2026EpistemicSybil}{55}{]}. Each
supporting record names the process or person that produced it. A record
without a producer, or without inputs unless it is a source observation,
counts as dependent and adds no root. Distinct roots are not
automatically independent: shared sources, producers and common causes
still need review. Supplement S4's draft schema gives the required
fields.

Independent quality review examines the cited records directly, weighs
contradicting evidence and alternative explanations, and checks
attribution and fit to the intended reader. It verifies claims of
detection validation against engineering evidence; detection development
and threat-model integration (responsibilities 16 and 17 in Supplement
S4) still perform that validation. Reviewers and checkers are subject to
the same untrusted-content safeguards as the agent.

When realistic volume forces reviewers to sample the cited records,
sampling is tiered by materiality. Claims about named entities are
checked in full, and every draft still requires human release. E8 tests
this policy under the planned workload.

Independence is a hypothesis, not a property guaranteed by a model name.
E6 and E8 measure whether review catches attacks and planted errors. E8
plants errors taken from the agent's own earlier runs alongside
independently written errors and reports the catch rate for each.
Separate sessions may share errors through common weights
{[}\protect\hyperlink{ref-Bara2026EpistemicSybil}{55}{]}, and different
models may share training data. A distinct method and organizational
role can provide separation, but the reviewer must still examine the
evidence.

Normal release follows quality review. The release decision records the
person, exact version hash, claims with the confidence stated for each,
scope, expiry and reviewer findings. The releaser sees supporting and
contradicting records beside each material claim. Release preserves the
assessment's uncertainty. Human release is itself exposed to automation
bias, in which people follow an automated aid's incorrect advice or miss
problems it fails to flag, and to the fall in detection when errors are
rare {[}\protect\hyperlink{ref-ParasuramanManzey2010}{56},
\protect\hyperlink{ref-Wolfe2005Rare}{57}{]}.

Urgent containment has a narrower route. A recommendation for a
reversible action, such as a payout pause or API-key restriction, goes
to an authorized on-call human with a decision target in minutes.
Quality review follows within a set period, and an authorized human or
accountable system decides whether to reverse the action. The draft does
not enter organizational memory through this exception. Existing
automatic controls remain in place.

Reversibility does not imply low impact. An API-key restriction may
interrupt a merchant's integration. Actions are classed by impact as
none, degraded service or outage; outage-class actions require
notification and a restoration path.

Human review also creates a capacity limit. With assumed values of 20
drafts a day and 30 minutes per draft, review alone requires 10 hours
daily. The system monitors release volume, returned drafts and the
oldest queued draft. Urgent containment takes priority when the queue is
saturated; expired drafts are discarded rather than released
automatically. E5 measures whether the agent's contribution justifies
this work, and E8 tests the gate under realistic volume and error
prevalence.

\hypertarget{protect-the-investigation-from-its-inputs}{%
\subsection{Protect the investigation from its
inputs}\label{protect-the-investigation-from-its-inputs}}

Businesses, customers and outside sources can supply content the agent
reads: storefront pages, descriptions, refund evidence, retrieved
documents and tool results. Indirect prompt injection uses such content
to introduce instructions. Research demonstrates these attacks through
external content and tool outputs
{[}\protect\hyperlink{ref-Greshake2023}{58},
\protect\hyperlink{ref-Debenedetti2024}{59}{]}.

The architecture treats this material as data with no authority over
instructions or tools. The gateway enforces task and tenant scope, least
privilege, the allowed tools, query limits and outbound-traffic
restrictions. These boundaries constrain access and action even if
hostile input influences reasoning.

\textbf{Approval binds to the operation actually reviewed.} It
identifies the exact payload, target, task, tenant, tool and
tool-manifest version, plus expiry and permitted reuse. At execution,
the gateway checks both the match and current permission. This addresses
altered operations and changed tool descriptions
{[}\protect\hyperlink{ref-Kumar2026}{60},
\protect\hyperlink{ref-InvariantLabs2025}{61}{]}. Approval of an
investigative operation cannot grant the agent authority over accounts,
payments or controls.

\textbf{Memory and outbound traffic have separate controls.} Proposed
lessons remain outside organizational memory until testing, quality
review and human release. This responds to documented memory-poisoning
routes {[}\protect\hyperlink{ref-Dash2026}{62}{]}. Traffic restrictions
respond to demonstrated exfiltration, in which injected instructions
passed user data to an attacker through search queries or links
{[}\protect\hyperlink{ref-Greshake2023}{58}{]}. Testing includes
adaptive attackers because fixed-test success can overstate resistance
{[}\protect\hyperlink{ref-Nasr2025}{63}{]}.

The threat model includes four attacker positions: an external abuser
controlling retrieved content; a merchant also supplying decisions and
refund evidence; a compromised operator; and an insider with approval or
release authority. Each plane's write rules must be evaluated against
all four. No control in this design stops the fourth position from
releasing a false claim or removing an agent-output or released-item
link, because its authority permits both. Against that position the
safeguards are the release record, the logged link removal, expiry and
re-testing of exonerating items, and correction (Section 7.3). They can
detect and limit such actions, but they do not prevent them.

One consequential attack produces a plausible but false benign
explanation. Human release alone does not defeat this attack.

Released benign assessments and exonerating lessons therefore expire on
a fixed date, or earlier when their evidence or cluster membership
changes. They are re-tested against independent counterexamples from
closed cases or reviewers who did not write the lesson. Labels,
assessments and decisions informed by them retain links so they can be
revisited after expiry or retraction.

\hypertarget{interpreting-feedback-after-protective-decisions}{%
\section{Interpreting feedback after protective
decisions}\label{interpreting-feedback-after-protective-decisions}}

This section answers RQ3. Section 4.2 distinguished three feedback
loops. This section concerns the two that leave the case: protective
feedback, which changes the platform and therefore the evidence
available afterwards (Sections 7.1 and 7.2), and organizational-learning
feedback, which carries conclusions and procedures into later cases
(Section 7.3).

\hypertarget{separate-delay-observability-and-report-truth}{%
\subsection{Separate delay, observability and report
truth}\label{separate-delay-observability-and-report-truth}}

A control changes both activity and the evidence later available about
it. Blocking a payment prevents that attempt from completing; holding a
payout may delay it. The absence of a later outcome therefore cannot be
interpreted without the decision that shaped its channel.

C1 keeps three questions separate (Figure 2): \textbf{Could a report
still arrive? Could this outcome occur under the decision taken? Is the
report true?} Their answers can depend on one another. A report may be
late and false, while a control may close one channel and leave another
open.

\begin{figure}
\centering
\includegraphics[width=\linewidth,height=0.5\textheight,keepaspectratio]{FIGURES/FIG2_FEEDBACK.pdf}
\Description{Figure 2: Decisions, delayed reports, outcomes made unobservable by controls and false reports}
\caption{\textbf{How decisions shape later evidence.} Delay concerns report
arrival; unobservability concerns outcomes a decision prevents; false
reports concern truth. Separate fields preserve these distinctions,
while evaluation must justify their joint assumptions. Corrected
intelligence and tested lessons return through review and release.}
\end{figure}

\hypertarget{delay-a-report-has-not-arrived-yet}{%
\subsubsection*{Delay: a report has not arrived
yet}\label{delay-a-report-has-not-arrived-yet}}
\addcontentsline{toc}{subsubsection}{Delay: a report has not arrived
yet}

A deployed payment-fraud system trains on user fraud reports, which
arrive after the transaction
{[}\protect\hyperlink{ref-Cao2019TitAnt}{5}{]}. The relevant reporting
window depends on the channel. A delay model can distinguish a pending
outcome from a negative after sufficient time, as in Chapelle's
conversion model {[}\protect\hyperlink{ref-Chapelle2014Delayed}{11}{]}.
Reading immature negatives as final can reverse a comparison: Stripe's
engineers report a fraud model trained on far more payments that looked
better at first, until, as more fraud reports arrived, its measured
performance fell below the earlier model's and it was rolled back
{[}\protect\hyperlink{ref-Drapeau2026ScalingRadar}{31}{]}. A channel
without a defined reporting window, such as a victim report, supplies no
such maturity bound.

The records preserve arrival time, channel clocks and the evaluation's
maturity choice. A late report may evaluate an earlier fixed prediction.
It must not enter that prediction's features or a training run whose
cutoff preceded its availability.

\hypertarget{observability-the-decision-changed-the-possible-outcome}{%
\subsubsection*{Observability: the decision changed the possible
outcome}\label{observability-the-decision-changed-the-possible-outcome}}
\addcontentsline{toc}{subsubsection}{Observability: the decision changed
the possible outcome}

A payment permanently declined before completion cannot generate a
chargeback on that attempt. It may still produce a complaint or incident
finding. A payout hold changes when funds move, not when a dispute can
arrive.

Stripe's dispute guide says card networks typically allow cardholders to
dispute within 120 days of the original payment, sometimes longer, and
that for a future event or service the window generally starts on the
event date {[}\protect\hyperlink{ref-Stripe2026Disputes}{49}{]}. Rules
vary by network and version. The records therefore distinguish
authorization, capture and processing times and use the clock-start
event specified by the applicable channel rule. Missing rules or
timestamps leave maturity unknown.

Capture may also gate possible-fraud reporting. Stripe's monitoring
guide says issuers must report possible fraud on captured payments, even
after refund, but are not required to report it on authorizations. It
adds that an authorization reversed before capture is not reported
{[}\protect\hyperlink{ref-Stripe2026MonitoringPrograms}{51}{]}. Stripe's
dispute guide adds that a refund processed as a reversal, which usually
happens within about 2 hours of capture, can also prevent a fraud report
{[}\protect\hyperlink{ref-Stripe2026Disputes}{49}{]}. Capture alone
therefore does not fix that channel's state. Thus a review ending in
reversal can close that channel while other incident reports remain
possible.

A control can change a channel's meaning without closing it. Stripe's
3-D Secure guide says liability for eligible fraud disputes can be
expected to shift to the issuer after successful authentication, but is
not guaranteed. It adds that an early fraud warning, a network report of
suspected fraud, may still arrive
{[}\protect\hyperlink{ref-Stripe2026ThreeDS}{64}{]}. Authentication
result, reported liability status and each channel's state must
therefore remain separate.

Payout timing is recorded against the relevant reporting windows because
funds may leave before reports arrive. The effect is recorded per
channel: closed, delayed, changed in meaning, none or unknown.

These channel properties come from card-network rules, and other rails
differ. On Stripe's bank-debit rails, a payer's dispute or return inside
the rail's no-questions window is honored without adjudication or appeal
{[}\protect\hyperlink{ref-Stripe2026ACHTerms}{52},
\protect\hyperlink{ref-Stripe2026SEPADebit}{65}{]}; Supplement S4.1
gives the windows and rules read. A report honored without question is
never adjudicated, so any won or lost status recorded for it says
nothing about its truth, which rests entirely on a separate assessment.
Each channel record therefore names its rail, the attributes its window
depends on and whether the channel adjudicates report truth.

\hypertarget{truth-the-report-may-be-wrong}{%
\subsubsection*{Truth: the report may be
wrong}\label{truth-the-report-may-be-wrong}}
\addcontentsline{toc}{subsubsection}{Truth: the report may be wrong}

Bad-faith disputes and fabricated refund evidence can make the adversary
a source of outcome labels
{[}\protect\hyperlink{ref-Dhama2026Causal}{8},
\protect\hyperlink{ref-Zhang2026Refund}{50}{]}. One evaluation interface
accepts the business's decision on each evaluation (allowed, restricted
or blocked) and says these reports train its models
{[}\protect\hyperlink{ref-Stripe2026CAE}{66}{]}. Such reports may also
carry that business's errors. Models trained on these labels predict
reported outcomes unless the reports' relationship to truth is
separately established.

Report type and source therefore remain attached to the label. Any
judgment that it is false is a separate assessment with its own
provenance. This preserves both the report and the uncertainty about its
truth.

\hypertarget{what-the-records-permit-an-evaluation-to-estimate}{%
\subsubsection*{What the records permit an evaluation to
estimate}\label{what-the-records-permit-an-evaluation-to-estimate}}
\addcontentsline{toc}{subsubsection}{What the records permit an
evaluation to estimate}

Selective-label methods and causal analysis require more than the
outcome field {[}\protect\hyperlink{ref-Lakkaraju2017Selective}{1},
\protect\hyperlink{ref-Kilbertus2020Fair}{2},
\protect\hyperlink{ref-HernanRobins2020}{4}{]}.

An observational evaluation must define its target population, the
decisions it compares and the outcome, then justify the assumptions
linking observed decisions to alternatives: consistency, conditional
exchangeability, overlap and no interference
{[}\protect\hyperlink{ref-HernanRobins2020}{4}{]}. The records support
examining assumptions; they do not validate an estimator or the bounds
in theoretical payment models
{[}\protect\hyperlink{ref-Dhama2026Causal}{8},
\protect\hyperlink{ref-Dhama2026Limits}{9}{]}. Supplement S7 defines the
four assumptions and illustrates the distinction with three target
quantities.

Without further assumptions, outcomes hidden by deterministic blocks
support ranges rather than a recovered comparison.

Interference is especially consequential when blocked actors move to
sibling accounts or other merchants. Links from each report to every
decision that may have shaped it make these dependencies visible; they
do not remove them.

Delay and false reports add assumptions of their own (Supplement S7).

These assumptions are joint. Allowed payments may receive chargebacks
while blocked items receive investigator labels or complaints, with
different delays and clocks. Separate fields make these differences
visible; they do not make the observation processes independent.

Obtaining missing comparisons through prospective exploration would mean
allowing some attempts the existing policy would block. That requires
separate ethical, human and operational authorization. It is neither
performed nor authorized here.

\hypertarget{attribution-and-resolution-change-measured-ratios}{%
\subsubsection*{Attribution and resolution change measured
ratios}\label{attribution-and-resolution-change-measured-ratios}}
\addcontentsline{toc}{subsubsection}{Attribution and resolution change
measured ratios}

Account structure affects where a report is recorded and which
transactions a ratio describes. FTC complaints allege pooling or
spreading activity to evade monitoring
{[}\protect\hyperlink{ref-FTC2020FirstData}{67},
\protect\hyperlink{ref-FTC2025Paddle}{68}{]}. These allegations motivate
transaction links and cluster-level measures, without establishing
erroneous attribution in a new case.

The same outcome definition, cohort and window must be retained when
comparing account and cluster ratios; Section 8.2 shows what pooling can
conceal.

Sales can also be inflated or moved. One FTC complaint alleges plans to
add stored-value card ``purchases'' to a merchant's sales, and continued
processing of a merchant's transactions through a linked account after
its account was closed {[}\protect\hyperlink{ref-FTC2019Allied}{69}{]}.
Payments known to be funded by instruments the business or platform
controls are therefore kept out of a ratio's denominator and counted
separately, and a closure's effect is measured across linked accounts.

Resolution also changes what becomes a chargeback. One FTC complaint
alleges refunding on pre-dispute alerts to avoid chargeback recording;
the resulting order was stipulated
{[}\protect\hyperlink{ref-FTC2025Paddle}{68}{]}. Stripe's dispute guide
says a business can proactively refund a payment flagged by an early
fraud warning to prevent a dispute, while advising against refunding
every flagged payment. It adds that such a refund does not affect the
fraud warning {[}\protect\hyperlink{ref-Stripe2026Disputes}{49}{]}. A
lower chargeback ratio after alert resolution therefore does not
establish less abuse.

The records retain each alert, its resolution and any later chargeback.
Transaction links allow a defined measure to count each affected
transaction once. Supplement Table S7.1 summarizes the resulting
requirements.

\hypertarget{measure-control-errors-and-user-costs-separately}{%
\subsection{Measure control errors and user costs
separately}\label{measure-control-errors-and-user-costs-separately}}

A known control decision can have unknown correctness. A study of five
commercial identity-verification products reported false negative rates
of 10.5 to 53.1 percent, the share of genuine participants' verification
attempts that were rejected
{[}\protect\hyperlink{ref-Fatima2024}{70}{]}. It did not report
fraud-detection results. Those cohort results cannot become
probabilities that an individual platform decision is wrong.

Denominators matter as much as percentages: Supplement S1 works through
one login-challenge study's rates
{[}\protect\hyperlink{ref-Doerfler2019}{71}{]}.

The specification links a control decision to a correctness review where
available and separately to applicable population error estimates. Each
estimate retains its version, cohort, outcome window, denominator,
uncertainty and scope. Evaluation weighs detection gain against
legitimate-user cost. Unknown, not-assessed and failed-open states
remain distinct from low risk
{[}\protect\hyperlink{ref-Stripe2026PAYG}{72}{]}.

Controls also need retesting as adversaries adapt. Published examples
include changed card-testing methods, aged accounts, banks changed after
payment disruption, and filtering of recognizable test transactions
{[}\protect\hyperlink{ref-Thomas2013}{53},
\protect\hyperlink{ref-Meltzer2025}{73},
\protect\hyperlink{ref-Tian2018}{74}{]}. Responsibilities 16 and 17
(Supplement S4) own this evaluation.

Test, replay and emulation records are marked when created and excluded
from production labels, intelligence and lessons unless their evaluation
use is explicitly authorized. The mark remains internal so it does not
help adversaries identify test traffic.

\hypertarget{correct-conclusions-and-test-lessons-before-reuse}{%
\subsection{Correct conclusions and test lessons before
reuse}\label{correct-conclusions-and-test-lessons-before-reuse}}

New evidence may require correcting or retracting released intelligence.
The correction links to the original, retains the new evidence's
provenance, repeats applicable review and release, and notifies
downstream users where warranted.

A reusable procedure follows a different path. A productive query, a
useful move between related entities, or a recurring benign explanation
becomes a proposed lesson. It enters a queue for prioritization and
independent tests against counterexamples. Lessons that exonerate expire
and are re-tested.

After testing, quality review and human release, a lesson may be reused
within a defined scope and version. It remains monitored and can be
withdrawn. This is controlled learning (responsibility 23 in Supplement
S4).

Memory-poisoning and experience-reuse studies show how planted or
mistaken entries can steer later work
{[}\protect\hyperlink{ref-Dash2026}{62},
\protect\hyperlink{ref-Xiong2026MemoryManagement}{75}{]}. A validated
lesson therefore guides procedure; it is never evidence that the
suspected activity occurred in the current case.

\hypertarget{two-hypothetical-investigations}{%
\section{Two hypothetical
investigations}\label{two-hypothetical-investigations}}

The following traces apply the specification to fixed hypothetical
observations. Each asks what happened, which explanations remain
plausible, and what further test would distinguish them. Their purpose
is to expose the consequences of missing context. Table 2 gives what the
recorded context changes. Supplement S6 gives each trace's reasoning and
outputs, the synthetic transaction records and reproducible arithmetic.

\hypertarget{trace-a-an-exposed-key-and-a-changed-payout-destination}{%
\subsection{Trace A: an exposed key and a changed payout
destination}\label{trace-a-an-exposed-key-and-a-changed-payout-destination}}

A repository scan finds a publicly visible API key, and a human asks
whether it enabled account compromise and payout redirection.

\textbf{Observations.} Nineteen minutes after publication, requests
arrive through a residential proxy range with a different client
fingerprint. The same burst creates a webhook endpoint, changes the
payout destination and accelerates the payout schedule; the trace
assumes the key permits these changes. The next day, delayed records
show that the same verified destination identifier was attached to two
other accounts before exposure.

\hypertarget{trace-b-linked-storefronts-and-dispute-activity}{%
\subsection{Trace B: linked storefronts and dispute
activity}\label{trace-b-linked-storefronts-and-dispute-activity}}

A risk assessment and a rejected application identify similar
storefronts, and a human asks whether newly onboarded merchant accounts
form an operation structured to evade dispute monitoring.

\textbf{Observations.} The hypothetical investigation links eleven
storefronts that share two templates and a descriptor pattern and settle
through four merchant accounts. A principal also appears on a closed
account and the rejected application. In a synthetic cohort of 100
completed transactions per account, every account has a chargeback ratio
of 2 percent once reports are linked to transactions. A separate
chargeback-or-alert-resolution measure, counting each transaction once,
is 5 percent pooled: 8 percent at accounts A and B and 2 percent at C
and D.

\hypertarget{what-retaining-the-context-changes}{%
\subsection{What retaining the context
changes}\label{what-retaining-the-context-changes}}

\begingroup\footnotesize

\begin{longtable}[]{@{}
  >{\raggedright\arraybackslash}p{(\columnwidth - 4\tabcolsep) * \real{0.3000}}
  >{\raggedright\arraybackslash}p{(\columnwidth - 4\tabcolsep) * \real{0.3400}}
  >{\raggedright\arraybackslash}p{(\columnwidth - 4\tabcolsep) * \real{0.3600}}@{}}
\caption{\textbf{Five conclusions that depend on recorded context.} Both
columns use the stipulated observations. The comparison is logical, not
a measured baseline. Supplement S4.4 gives four C2 and C3 dependencies
as pairs of record histories that permit different conclusions but are
identical once one group of fields is removed.}\tabularnewline
\toprule\noalign{}
\begin{minipage}[b]{\linewidth}\raggedright
Question and fixed observations
\end{minipage} & \begin{minipage}[b]{\linewidth}\raggedright
If context is omitted
\end{minipage} & \begin{minipage}[b]{\linewidth}\raggedright
If context is retained
\end{minipage} \\
\midrule\noalign{}
\endfirsthead
\toprule\noalign{}
\begin{minipage}[b]{\linewidth}\raggedright
Question and fixed observations
\end{minipage} & \begin{minipage}[b]{\linewidth}\raggedright
If context is omitted
\end{minipage} & \begin{minipage}[b]{\linewidth}\raggedright
If context is retained
\end{minipage} \\
\midrule\noalign{}
\endhead
\bottomrule\noalign{}
\endlastfoot
Did misuse begin after exposure? A seven-day query returns two hours of
earlier history; changed requests appear at +19 minutes. & The lack of
an earlier returned match does not distinguish absent activity from
missing history. & The baseline is partial. The change is observed, but
the onset of misuse remains inconclusive. \\
Could the first assessment use destination reuse? The link was effective
before exposure, learned at +1 day; the assessment was frozen at +30
minutes. & Effective-date ordering can introduce information unavailable
at the decision. & Learned time excludes the link from that assessment.
It may inform a later investigation; later labels may still evaluate the
fixed assessment. \\
Can a blocked payout receive a negative outcome label? The transfer
never completes. & The absence of a return or victim report can be
mistaken for a benign outcome. & The return-or-recall channel is closed,
not negative. Victim-report and incident-finding channels remain open.
The blocked case remains in any target cohort covering all attempts. \\
Which transactions do reports concern? Four cohorts contain 400
transactions, eight with chargebacks; D001 and D002 are recorded at
account-X. Two of fourteen alert resolutions overlap. &
Recording-account counts show D at 0/100. Adding report flags gives
22/400 and changes the outcome definition. & Transaction links restore D
to 2/100 and remove double counting. Chargebacks are 8/400 = 2\%; the
broader measure is 20/400 = 5\%, with A and B at 8/100 and C and D at
2/100. Neither establishes fraud. \\
Does an earlier released assessment corroborate the cluster finding? It
concerns the same descriptor pattern. & The conclusion can become
apparent independent support for itself. & The assessment supplies a
hypothesis or context. It adds no root observation and no independent
corroboration. \\
\end{longtable}

\endgroup

Each row of Table 2 isolates one requirement. The first two rows of
Table 1 are pairs of histories in which two requirements hold and only
the third shows whether an absence is a negative (Supplement S4.4).

Two rows of Table 2 add explicit hypothetical conditions: Trace A's
payout is permanently blocked, and Trace B retrieves an earlier
assessment of the same descriptor pattern. Omission does not force an
error; it removes information needed to rule the error out.

That context costs storage, query work, controlled cross-account access
and investigator time. E1 tests the context in the first two rows of
Table 2, E9 the channel states in the third, E3 the report linkage in
the fourth, and E5 the human work the agent creates. E8's planted drafts
that rest a material claim only on a released item test the rule in the
fifth row. Section 8.4 estimates the records' storage and query work on
synthetic data.

\hypertarget{where-to-start-and-what-it-costs}{%
\subsection{Where to start and what it
costs}\label{where-to-start-and-what-it-costs}}

A platform that already logs decisions, payments and disputes can start
with the records behind Table 2: each decision's effect on each outcome
channel and each channel's state for each payment (C1), a gateway record
of what each query covered and when each fact became known (C2), and a
release step before an agent conclusion enters reusable memory (C3).
Stripe already exposes parts of the first two (Section 2.4).

To estimate C1's cost, one million synthetic payments with two outcome
channels each were built from the schema's valid conformance case
(Supplement S4.3). Each payment needed about four records; its decision
and attribution records may already exist, and its two channel-state
records are new. The records took 1.7 GB per million payments as JSON
and 2.6 GB in a single-file SQLite store, and reading one payment's
channel state and reports took a median of 17 microseconds. A query
run's coverage record took about 900 bytes, and checking an absence
claim against its runs took 38 microseconds (Supplement S4.4). These
figures describe one synthetic file on one machine, not a platform's
cost. Reviewer time is not measured.

\hypertarget{synthetic-replays-of-two-alleged-patterns}{%
\subsection{Synthetic replays of two alleged
patterns}\label{synthetic-replays-of-two-alleged-patterns}}

Two FTC complaints were rebuilt as synthetic outcome records from the
patterns and rates they allege, with illustrative monitoring rules and
assumed volumes and delays (Supplement S10)
{[}\protect\hyperlink{ref-FTC2020FirstData}{67},
\protect\hyperlink{ref-FTC2019Allied}{69}{]}. Complaint content is
allegation; the replays test record design, not detection, and whether a
view flags follows from the chosen parameters. The records have limits:
padded sales look like clean payments, a record rewritten in place
passes every rule, so the store must prevent edits, and accounts with no
recorded shared attribute cannot be linked.

Network programs compute monthly ratios by the month a report is
received, apply minimum counts, and in one program count early fraud
warnings and exclude disputes resolved through pre-dispute products
{[}\protect\hyperlink{ref-Stripe2026MonitoringPrograms}{51}{]};
Supplement S6 gives the rules as read and where Visa's own overview
differs {[}\protect\hyperlink{ref-Visa2025VAMP}{76}{]}. Splitting
activity can also keep accounts below a count threshold even when their
ratios are identical. Detecting that pattern requires cluster-level
counts, not merely a pooled cohort ratio.

\begin{figure}
\centering
\includegraphics[width=\linewidth,height=0.5\textheight,keepaspectratio]{FIGURES/FIG3_REPLAY_SWEEPS.pdf}
\Description{Figure 3: Share of runs flagged by each view as volume, report rate and payments a month change}
\caption{\textbf{Where each view stops flagging.} The share of 500 synthetic runs
in which a view flags at least once. (A, B) First Data replay: eight
related accounts, at a 4 percent report rate in A and 350 captures each
in B, and 40 merchants assumed legitimate, at 300 each; the dotted line
is the 1.5 percent ratio. Any flag of the legitimate merchants is false
by assumption. (C) Allied Wallet replay: one account at the complaint's
mean fraud-report rate of 3.14 percent or at 0.3 percent. Supplement S10
gives the assumptions.}
\end{figure}

In the first replay, the account view flags only when each account alone
reaches the floor, at about 2,500 captures a period (Figure 3A). A view
over recorded links flags the group from about 350. Below 300 the whole
group stays under the floor, and neither view flags in more than 1
percent of runs. The linked view also flags legitimate groups (Figure
3B). The 40 merchants have enough volume that the ratio, not the floor,
decides. Between report rates of 1.75 and 3 percent, the linked view
flags them in at least 98 percent of runs and the smaller related group
in at most 2.2 percent. A linked view belongs beside the account view,
not in its place.

In the second replay, the chargeback rule needs 100 chargebacks a month.
With a fifth of reported payments refunded before a dispute, one account
needs about 4,000 payments a month to reach it; the modeled accounts had
185 (Figure 3C). The complaint says Visa flags fraud reports at either
\$75,000 a month or 1 percent of sales by value
{[}\protect\hyperlink{ref-FTC2019Allied}{69}{]}. Counted that way, fraud
reports flag in at least 98 percent of runs at every volume swept, from
25 a month; that rule needs the separate fraud-report channel and each
report's amount. With no count floor, it also flags an account at an
assumed legitimate rate of 0.3 percent in 39 to 54 percent of runs below
200 payments a month, and in none from 2,000. Supplement S10 also
replays a successor account, a switched descriptor, refunds before
disputes, traffic moved after closures and reports linked only by card
number and amount.

\hypertarget{evaluating-the-specification}{%
\section{Evaluating the
specification}\label{evaluating-the-specification}}

\hypertarget{questions-comparisons-and-failure-criteria}{%
\subsection{Questions, comparisons and failure
criteria}\label{questions-comparisons-and-failure-criteria}}

The specification is worthwhile only if its added records and safeguards
improve investigation enough to justify their cost. The evaluation plan
separates the records from the agent: C1 and C2 do not need one, E1, E3
and E9 test them without it, and E5 tests the agent's benefit. A
platform can start within a quarter with E9(a) and an audit of closed
cases that measures how often each targeted failure occurs (Supplement
S7.1).

Supplement S7.1 gives each hypothesis, design, primary measure and
proposed failure criteria. E1, E3, E8 and E9 are confirmatory; the
others are exploratory. The tests address the three requirements and the
agent's added value; none tests the signal inventory or the
responsibilities as a whole. E3 tests the report-to-transaction links on
which C1's outcome records depend, together with cluster-level
monitoring. Its pass rule compares linked clusters with accounts as
their reports are recorded, so a lead from clustering alone can pass it;
only the third arm, which applies the rule to each account after
linking, shows whether linkage contributes. E9 exercises C1's links from
outcomes to control decisions, and its separation of late from
unobservable outcomes, only jointly through the channel-state reading.
No confirmatory evaluation tests the separate field for false reports or
the decision links' other uses, such as estimating outcomes under other
decisions. The exploratory E2 bears on the decision links and, through
its sensitivity analysis, on late and false reports; E4 bears on the
decision links. Passing E1, E3, E8 and E9 would therefore not establish
the value of each of those links and fields on its own. E8 tests C3's
release gate. The exploratory E6 checks that required links back to
agent outputs exist. No evaluation tests the urgent-containment route in
Section 6.4. The rate at which later review reverses its actions, and
the service those actions interrupt, remain unmeasured. The questions
below state what each evaluation is intended to resolve.

\textbf{E1. Coverage and learned time:} Do these records reduce
unsupported conclusions in cases with incomplete history or links
learned late, without reducing the correct, supported conclusions
investigators reach in those cases or in cases with complete history?

\textbf{E2. Unobserved outcomes:} For a defined outcome and population,
which identifying assumptions (the assumptions needed to estimate
outcomes that could not be observed) are defensible? Where is overlap
(Section 7.1) absent? How sensitive are estimates to late, mistaken or
deliberately false reports?

\textbf{E3. Report linkage and cluster monitoring:} Does linking reports
to transactions and applying the network's monthly rule, its ratio, its
minimum count and any minimum value, to each cluster flag evasion
earlier than account-level monitoring? A broader measure, such as
Section 8.2's chargeback-or-alert-resolution measure, needs a separate
evaluation.

\textbf{E4. Measuring user harm:} Can the records support a rate and
interval for harm to legitimate users after a control change?

\textbf{E5. Agent benefit:} When all have the same information, does the
agent with human review match or improve on the confirmed findings per
reviewer hour of well-equipped human investigators, without materially
more errors than either them or suitable decision rules that use no
agent?

\textbf{E6. Safeguards:} How well do the agent's safeguards, covering
prompt injection, approvals, its memory and the data it sends out, hold
up against attackers who adapt to them? How much do they reduce the
number of investigations completed?

\textbf{E7. Learning:} How quickly do useful lessons become available
for reuse? How do testing against counterexamples, scoped promotion and
withdrawal affect negative transfer (a reused lesson degrading later
investigations)? How do they affect poisoning (attackers deliberately
causing false lessons)?

\textbf{E8. Release gate:} Does the release gate catch planted wrong
drafts under realistic volume and error prevalence, when reviewers and
releasers do not know which drafts are planted, and how much reviewer
time does it require? Some planted drafts rest a material claim only on
a released item or on records linked to one, which tests C3's rule that
such support adds no independent corroboration. Its criteria set no
limit on correct drafts that the gate returns or lets expire. Passing E8
would show that the gate catches errors, not that it releases correct
work at an acceptable rate; E5 bears on that cost.

\textbf{E9. Channel states:} Do C1's channel states reduce the absolute
error of evaluation results against a mature holdout reference, compared
with the best reading without them (each rail's standard window, blocked
attempts excluded)? The target is the selected metric on matured
reported outcomes after a sampled review of report truth, not underlying
fraud that may never be reported. E9 reports the share of items that C1
reclassifies but does not count it as support, because that share
follows partly from C1's definitions and the cohort; support requires a
reduction in reference error above the prespecified negligible
improvement.

Supplement S8 gives a practitioner study, fixed in advance, of whether
the records change practitioners' conclusions.

Supplement S7 gives the conditions and motivating sources for E1 to E7
and E9, and S7.1 states the tests for all nine. Four rules in S7.1 apply
to every evaluation: an as-of cutoff on the information each method may
use, outcomes decided blind under rules set in advance, a stated
treatment of labels shaped by control decisions, and marking of test
records. Of the S7.1 tests, E5 and E8 measure reviewer time, and E5 also
compares cost across arms that see the same records. None measures the
storage, query work or controlled cross-account access that the added
records require. S7 adds collection cost for E1's broader timing study
and the safeguards' effect on investigations completed for E6. Passing
the S7.1 tests would therefore not show that the added records justify
their collection cost. E1, E5 and E8 also measure the work of
identifiable staff on real case records, and E8's undisclosed planted
drafts require a deception protocol. Each needs ethics and privacy
review before it starts.

\hypertarget{what-the-evidence-can-support}{%
\subsection{What the evidence can
support}\label{what-the-evidence-can-support}}

\begingroup\footnotesize

\begin{longtable}[]{@{}
  >{\raggedright\arraybackslash}p{(\columnwidth - 4\tabcolsep) * \real{0.3000}}
  >{\raggedright\arraybackslash}p{(\columnwidth - 4\tabcolsep) * \real{0.3400}}
  >{\raggedright\arraybackslash}p{(\columnwidth - 4\tabcolsep) * \real{0.3600}}@{}}
\caption{\textbf{Each claim, its evidence and what that evidence does
not show.} Supplement S9.6 gives further limits.}\tabularnewline
\toprule\noalign{}
\begin{minipage}[b]{\linewidth}\raggedright
Claim
\end{minipage} & \begin{minipage}[b]{\linewidth}\raggedright
Evidence
\end{minipage} & \begin{minipage}[b]{\linewidth}\raggedright
Not shown
\end{minipage} \\
\midrule\noalign{}
\endfirsthead
\toprule\noalign{}
\begin{minipage}[b]{\linewidth}\raggedright
Claim
\end{minipage} & \begin{minipage}[b]{\linewidth}\raggedright
Evidence
\end{minipage} & \begin{minipage}[b]{\linewidth}\raggedright
Not shown
\end{minipage} \\
\midrule\noalign{}
\endhead
\bottomrule\noalign{}
\endlastfoot
Each field group changes what a history supports & Necessity pairs;
failing cases (S4.2, S4.4) & That the records improve investigations \\
No two requirements replace the third & Two cross-requirement pairs
(8.3, S4.4) & That they cover every interaction \\
Software can check the rules & Two checkers; 15 adversarial histories
(S4.4) & Every evasion; four stated limits remain \\
The records are cheap to store and check & Synthetic measurement (8.4) &
The cost of producing them on a platform \\
Missing context changes conclusions & Two hypothetical traces (8.1 to
8.3) & How often; the practitioner study (S8) has not run \\
The records expose two alleged patterns & Synthetic replays and sweeps
(8.5, S10) & Detection on real data \\
The signals and responsibilities cover the mechanisms & The author's
mapping; a model as second rater (3, S3) & Completeness; agreement
between human coders \\
No source states C1 to C3 together & Two searches (2, S9.6) & That none
exists \\
C1's channel fields reflect network rules & Stripe's descriptions (7.1)
& The networks' own rules; other rails reviewed less fully \\
\end{longtable}

\endgroup

A crosswalk to MITRE's Fight Fraud Framework finds that 21 of its 105
in-scope techniques have no retained mechanism; fourteen of them are
covered only by candidates outside the signal map (Supplement S2.2)
{[}\protect\hyperlink{ref-MITRE2026F3}{77}{]}. The prior-work comparison
supports the specified distinctions, not priority for their individual
components.

Evidence strength varies by premise. Several of the closest precedents
in Section 2 are 2026 preprints, and Dhama's payment models are
theoretical preprints {[}\protect\hyperlink{ref-Dhama2026Causal}{8},
\protect\hyperlink{ref-Dhama2026Limits}{9}{]}. The general
selective-label problem has broader support
{[}\protect\hyperlink{ref-Lakkaraju2017Selective}{1},
\protect\hyperlink{ref-Kilbertus2020Fair}{2}{]}, and the specification
does not depend on those preprints' estimators or bounds being correct.
Estimating a decision's effect still requires identifying assumptions
{[}\protect\hyperlink{ref-HernanRobins2020}{4}{]}. Enforcement
complaints motivate alleged mechanisms; vendor accounts support
attributed observations within their stated settings.

\hypertarget{conclusion}{%
\section{Conclusion}\label{conclusion}}

An investigation across payment products needs more than joined records.
It needs to preserve why an outcome could be observed, what information
was available to a query, and whether apparent support comes from an
earlier conclusion.

This paper combines those requirements in one specification grounded in
public evidence. The specification proposes answers to the three
research questions. For RQ1, records keep their meaning across products
through the seven record classes, the store contract and C2's coverage
and learned-time records, which bound what an investigation could have
established. Each analytical capability carries the limit on what its
results can show. For RQ2, the agent tests explanations without
operational authority. It works through governed queries whose coverage
it cannot alter, and its drafts pass a checker, independent quality
review and the human release of C3. Release authorizes reuse of a
conclusion from organizational memory, while the conclusion keeps its
status as context and its dependence on earlier agent outputs. For RQ3,
feedback informs later work without concealing control effects or
recycling conclusions. It rests on the decision-linked outcome channels
of C1, which keep missing reports from silently becoming negatives, and
on lessons that are tested, reviewed and released for a defined scope
before reuse, monitored during it and, when they exonerate, re-tested at
expiry. Decisions about accounts, payments and controls stay with
authorized humans or accountable operational systems.

For a platform that already logs decisions, queries and disputes, the
specification states what those records and their joins must preserve.
For designers and auditors of investigation systems, it places an
agent's access, its evidential claims and its reuse of conclusions in
one inspectable design.

The next test is operational value. Evaluations must determine whether
these records improve conclusions, whether the release gate works under
realistic demand, and whether the agent saves more work than review
consumes. The specification supplies a concrete basis for those tests
without presuming their result.

\hypertarget{author-and-publication-declarations}{%
\section*{Author and publication
declarations}\label{author-and-publication-declarations}}
\addcontentsline{toc}{section}{Author and publication declarations}

\textbf{Use of generative AI.} The author used Anthropic Claude and
OpenAI ChatGPT/Codex in the research, as Section 3 describes, and for
drafting and revision. The author takes responsibility for the content.

\textbf{Data and code availability.} The accompanying supplement
provides the signal inventory, mappings, entity and interface records,
source locators, schema pins and synthetic worked-example data.
Reproduction scripts check the included counts, lineage and arithmetic
and rebuild the documents. No private platform data, implementation of
the proposed architecture or empirical system-evaluation code is
supplied. External source texts remain with their publishers or authors.

\textbf{Funding.} This work received no funding.

\textbf{License.} Creative Commons Attribution 4.0 International (CC BY
4.0) applies to the author's text, figures and supplement tables. The
Stripe schema-derived components of two supplement tables remain under
the MIT License, as THIRD\_PARTY\_NOTICES.md states.

\textbf{Ethics and confidentiality.} No human participants were involved
in the results reported here. The paper uses only public sources, and no
confidential information from any current or former employer was used.
The practitioner study in Supplement S8 follows ACM's policy on research
involving human participants.

\hypertarget{references}{%
\section*{References}\label{references}}
\addcontentsline{toc}{section}{References}

\hypertarget{refs}{}
\begin{CSLReferences}{1}{1}
\leavevmode\vadjust pre{\hypertarget{ref-Lakkaraju2017Selective}{}}%
{[}1{]} H. Lakkaraju, J. Kleinberg, J. Leskovec, J. Ludwig, and S.
Mullainathan. 2017. The Selective Labels Problem: Evaluating Algorithmic
Predictions in the Presence of Unobservables. \emph{Proceedings of the
23rd ACM SIGKDD International Conference on Knowledge Discovery and Data
Mining}, pp. 275--284. \url{https://doi.org/10.1145/3097983.3098066}.

\leavevmode\vadjust pre{\hypertarget{ref-Kilbertus2020Fair}{}}%
{[}2{]} N. Kilbertus, M. Gomez-Rodriguez, B. Schölkopf, K. Muandet, and
I. Valera. 2020. Fair Decisions Despite Imperfect Predictions.
\emph{Proceedings of the Twenty Third International Conference on
Artificial Intelligence and Statistics}, 108, pp. 277--287.
\url{https://proceedings.mlr.press/v108/kilbertus20a.html}.

\leavevmode\vadjust pre{\hypertarget{ref-Bottou2013Counterfactual}{}}%
{[}3{]} L. Bottou, J. Peters, J. Quiñonero-Candela, D. X. Charles, D. M.
Chickering, E. Portugaly, D. Ray, P. Simard, and E. Snelson. 2013.
Counterfactual Reasoning and Learning Systems: The Example of
Computational Advertising. \emph{Journal of Machine Learning Research},
14(101), pp. 3207--3260.
\url{https://jmlr.org/papers/v14/bottou13a.html}.

\leavevmode\vadjust pre{\hypertarget{ref-HernanRobins2020}{}}%
{[}4{]} M. A. Hernán and J. M. Robins. 2020. Causal Inference: What If.
Boca Raton, Chapman \& Hall/CRC. Chapters cited: 1, 2, 3, 7 and 19.
Cited version: author-hosted revision (19 August 2026).
\url{https://miguelhernan.org/whatifbook}.

\leavevmode\vadjust pre{\hypertarget{ref-Cao2019TitAnt}{}}%
{[}5{]} S. Cao, X. Yang, C. Chen, J. Zhou, X. Li, and Y. Qi. 2019.
TitAnt: Online Real-time Transaction Fraud Detection in Ant Financial.
\emph{Proceedings of the VLDB Endowment}, 12(12), pp. 2082--2093.
\url{https://doi.org/10.14778/3352063.3352126}.

\leavevmode\vadjust pre{\hypertarget{ref-DalPozzolo2017Fraud}{}}%
{[}6{]} A. Dal Pozzolo, G. Boracchi, O. Caelen, C. Alippi, and G.
Bontempi. 2018. Credit Card Fraud Detection: A Realistic Modeling and a
Novel Learning Strategy. \emph{IEEE Transactions on Neural Networks and
Learning Systems}, 29(8), pp. 3784--3797.
\url{https://doi.org/10.1109/TNNLS.2017.2736643}.

\leavevmode\vadjust pre{\hypertarget{ref-Tax2021Fraud}{}}%
{[}7{]} N. Tax, K. J. de Vries, M. de Jong, N. Dosoula, B. van den
Akker, J. Smith, O. Thuong, and L. Bernardi. 2021. Machine Learning for
Fraud Detection in E-Commerce: A Research Agenda. arXiv preprint
2107.01979v1. \url{https://arxiv.org/abs/2107.01979v1}.

\leavevmode\vadjust pre{\hypertarget{ref-Dhama2026Causal}{}}%
{[}8{]} G. Dhama. 2026. Causal Label Recovery in Payment Networks. arXiv
preprint 2605.29272v1. \url{https://arxiv.org/abs/2605.29272v1}.

\leavevmode\vadjust pre{\hypertarget{ref-Dhama2026Limits}{}}%
{[}9{]} G. Dhama. 2026. The Fundamental Limits of Fraud Detection in
Card Payment Networks: An Information-Theoretic Theory of Delayed,
Censored, Corrupted, and Counterfactual Feedback. arXiv preprint
2605.27557v2. \url{https://arxiv.org/abs/2605.27557v2}.

\leavevmode\vadjust pre{\hypertarget{ref-HeLoh2026Grab}{}}%
{[}10{]} M. He and J. L. Loh. 2026. Combating Organized Platform Abuse:
Amplifying Weak Risk Signals with Structural Information. arXiv preprint
2605.07383v2. \url{https://arxiv.org/abs/2605.07383v2}.

\leavevmode\vadjust pre{\hypertarget{ref-Chapelle2014Delayed}{}}%
{[}11{]} O. Chapelle. 2014. Modeling Delayed Feedback in Display
Advertising. \emph{Proceedings of the 20th ACM SIGKDD International
Conference on Knowledge Discovery and Data Mining}, pp. 1097--1105.
\url{https://doi.org/10.1145/2623330.2623634}.

\leavevmode\vadjust pre{\hypertarget{ref-Solozobov2026DecisionTrace}{}}%
{[}12{]} O. Solozobov. 2026. Decision Trace Schema for Governance
Evidence in Real-Time Risk Systems. arXiv preprint 2604.09296v1.
\url{https://arxiv.org/abs/2604.09296v1}.

\leavevmode\vadjust pre{\hypertarget{ref-Wu2026EvidenceCoverage}{}}%
{[}13{]} M. Wu, L. Wang, Y. Zhang, and J. Deng. 2026. Evidence Coverage
for Intent-Bound Execution: Scope, Obligations, and Cutoff Reasoning.
arXiv preprint 2610.01386v1. \url{https://arxiv.org/abs/2610.01386v1}.

\leavevmode\vadjust pre{\hypertarget{ref-Snodgrass1986Temporal}{}}%
{[}14{]} R. Snodgrass and I. Ahn. 1986. Temporal Databases.
\emph{Computer}, 19(9), pp. 35--42.
\url{https://doi.org/10.1109/MC.1986.1663327}.

\leavevmode\vadjust pre{\hypertarget{ref-Wang2026TOKI}{}}%
{[}15{]} Z. Wang. 2026. TOKI: A Bitemporal Operator Algebra for
Contradiction Resolution in LLM-Agent Persistent Memory. arXiv preprint
2606.06240v1. \url{https://arxiv.org/abs/2606.06240v1}.

\leavevmode\vadjust pre{\hypertarget{ref-Rasmussen2025Zep}{}}%
{[}16{]} P. Rasmussen, P. Paliychuk, T. Beauvais, J. Ryan, and D.
Chalef. 2025. Zep: A Temporal Knowledge Graph Architecture for Agent
Memory. arXiv preprint 2501.13956v1.
\url{https://arxiv.org/abs/2501.13956v1}.

\leavevmode\vadjust pre{\hypertarget{ref-Dan2026Anamnesis}{}}%
{[}17{]} A. Dan. 2026. Anamnesis: An Authority-Preserving Runtime
Contract for Persistent AI Memory. Zenodo preprint, Version 1.0, 16
August 2026. \url{https://zenodo.org/records/22118090}.

\leavevmode\vadjust pre{\hypertarget{ref-Razniewski2011Completeness}{}}%
{[}18{]} S. Razniewski and W. Nutt. 2011. Completeness of Queries over
Incomplete Databases. \emph{Proceedings of the VLDB Endowment}, 4(11),
pp. 749--760. \url{https://doi.org/10.14778/3402707.3402715}.

\leavevmode\vadjust pre{\hypertarget{ref-W3C2013PROVDM}{}}%
{[}19{]} World Wide Web Consortium. 2013. PROV-DM: The PROV data model.
W3C Recommendation.
\url{https://www.w3.org/TR/2013/REC-prov-dm-20130430/}.

\leavevmode\vadjust pre{\hypertarget{ref-Singh2019DecisionProvenance}{}}%
{[}20{]} J. Singh, J. Cobbe, and C. Norval. 2019. Decision Provenance:
Harnessing Data Flow for Accountable Systems. \emph{IEEE Access}, 7, pp.
6562--6574. \url{https://doi.org/10.1109/ACCESS.2018.2887201}.

\leavevmode\vadjust pre{\hypertarget{ref-OASIS2021STIX}{}}%
{[}21{]} OASIS. 2021. STIX Version 2.1. OASIS Standard.
\url{https://docs.oasis-open.org/cti/stix/v2.1/os/stix-v2.1-os.html}.

\leavevmode\vadjust pre{\hypertarget{ref-ODNI2015ICD203}{}}%
{[}22{]} Office of the Director of National Intelligence. 2015.
Intelligence Community Directive 203: Analytic Standards. Directive. As
amended by technical amendments through 21 January 2022.
\url{https://archive.dni.gov/files/documents/ICD/ICD-203.pdf}.

\leavevmode\vadjust pre{\hypertarget{ref-Dai2026Hearsay}{}}%
{[}23{]} J. Dai, Z. Yang, P. Shao, K. Ng, Z. Liu, C. Ju, Y. He, and B.
Hu. 2026. Hearsay: Can an Auditor Trust the Record a Deployed Agent
Harness Writes? arXiv preprint 2609.32495v1.
\url{https://arxiv.org/abs/2609.32495v1}.

\leavevmode\vadjust pre{\hypertarget{ref-Doyle1979TMS}{}}%
{[}24{]} J. Doyle. 1979. A Truth Maintenance System. \emph{Artificial
Intelligence}, 12(3), pp. 231--272.
\url{https://doi.org/10.1016/0004-3702(79)90008-0}.

\leavevmode\vadjust pre{\hypertarget{ref-Malheiro1994BeliefRevision}{}}%
{[}25{]} B. Malheiro, N. R. Jennings, and E. Oliveira. 1994. Belief
Revision in Multi-Agent Systems. \emph{Proceedings of the 11th European
Conference on Artificial Intelligence}, pp. 294--298.
\url{https://eprints.soton.ac.uk/252143}.

\leavevmode\vadjust pre{\hypertarget{ref-Sumers2024CoALA}{}}%
{[}26{]} T. R. Sumers, S. Yao, K. Narasimhan, and T. L. Griffiths. 2024.
Cognitive Architectures for Language Agents. \emph{Transactions on
Machine Learning Research}. \url{https://arxiv.org/abs/2309.02427v3}.

\leavevmode\vadjust pre{\hypertarget{ref-Wang2026AgentTraces}{}}%
{[}27{]} Y. Wang, J. Zhang, Z. Wu, T. Cai, Z. Liu, Z. Sun, M. Dong, M.
Zheng, Y. Duan, X. Yin, and Y. Zhu. 2026. From Agent Traces to Trust: A
Survey of Evidence Tracing and Execution Provenance in LLM Agents. arXiv
preprint 2606.04990v5. \url{https://arxiv.org/abs/2606.04990v5}.

\leavevmode\vadjust pre{\hypertarget{ref-Qi2026GovMem}{}}%
{[}28{]} Y. Qi, X. Xu, and Y. Li. 2026. When Not to Write Memory:
Governing False Promotion from Correlated Agent Traces. arXiv preprint
2607.02579v1. \url{https://arxiv.org/abs/2607.02579v1}.

\leavevmode\vadjust pre{\hypertarget{ref-Brown2026Quipu}{}}%
{[}29{]} S. Brown. 2026. Quipu: A Governed Bitemporal Knowledge Graph
Store. arXiv preprint 2608.16813v1.
\url{https://arxiv.org/abs/2608.16813v1}.

\leavevmode\vadjust pre{\hypertarget{ref-Akhtar2026Contestable}{}}%
{[}30{]} H. Akhtar and B. Freeman. 2026. Contestable Agents for Data
Governance: Recording Human Stewardship of AI-Generated Metadata Across
Structured and Unstructured Systems. arXiv preprint 2610.31553v1.
\url{https://arxiv.org/abs/2610.31553v1}.

\leavevmode\vadjust pre{\hypertarget{ref-Drapeau2026ScalingRadar}{}}%
{[}31{]} R. Drapeau and S. Samant. 2026. Scaling Radar: Fighting fraud
with agent-assisted machine learning. Stripe developer blog.
\url{https://stripe.dev/blog/scaling-radar-fighting-fraud-with-agent-assisted-machine-learning}.
Accessed 26 October 2026.

\leavevmode\vadjust pre{\hypertarget{ref-Evron2026AdversaryBrains}{}}%
{[}32{]} G. Evron. 2026. Adversary Brains: Hypothesis-Based Logical
Signatures for Threat Actor Detection with LLMs. arXiv preprint
2610.35624v1. \url{https://arxiv.org/abs/2610.35624v1}.

\leavevmode\vadjust pre{\hypertarget{ref-Liu2018GEM}{}}%
{[}33{]} Z. Liu, C. Chen, X. Yang, J. Zhou, X. Li, and L. Song. 2018.
Heterogeneous Graph Neural Networks for Malicious Account Detection.
\emph{Proceedings of the 27th ACM International Conference on
Information and Knowledge Management}.
\url{https://doi.org/10.1145/3269206.3272010}.

\leavevmode\vadjust pre{\hypertarget{ref-Stripe2026QueryDisputes}{}}%
{[}34{]} Stripe. 2026. Query disputes and fraud data. Stripe
documentation.
\url{https://docs.stripe.com/data/query-disputes-and-fraud-data}.
Accessed 26 October 2026.

\leavevmode\vadjust pre{\hypertarget{ref-Stripe2026DisputeAnalytics}{}}%
{[}35{]} Stripe. 2026. Disputes analytics. Stripe documentation.
\url{https://docs.stripe.com/payments/analytics/disputes}. Accessed 26
October 2026.

\leavevmode\vadjust pre{\hypertarget{ref-Stripe2026SigmaRun}{}}%
{[}36{]} Stripe. 2026. The Scheduled Query object. Stripe documentation.
\url{https://docs.stripe.com/api/sigma/scheduled_queries/object}.
Accessed 26 October 2026.

\leavevmode\vadjust pre{\hypertarget{ref-Adyen2026RiskMonitoring}{}}%
{[}37{]} Adyen. 2026. Monitor risk performance. Adyen documentation.
\url{https://docs.adyen.com/risk-management/dispute-and-fraud-monitoring}.
Accessed 26 October 2026.

\leavevmode\vadjust pre{\hypertarget{ref-Solozobov2026Synthesis}{}}%
{[}38{]} O. Solozobov. 2026. Governed Auditable Decisioning Under
Uncertainty: Synthesis and Agentic Extension. arXiv preprint
2604.19112v2. Cited version: arXiv 2604.19112v2 (2026).
\url{https://arxiv.org/abs/2604.19112v2}.

\leavevmode\vadjust pre{\hypertarget{ref-MindForge2026Operationalisation}{}}%
{[}39{]} MindForge Consortium. 2026. AI Risk Management:
Operationalisation Handbook. Handbook, Monetary Authority of Singapore.
\url{https://www.mas.gov.sg/-/media/mas-media-library/schemes-and-initiatives/ftig/project-mindforge/mindforge-ai-risk-management-operationalisation-handbook.pdf}.

\leavevmode\vadjust pre{\hypertarget{ref-MAS2026SAFR}{}}%
{[}40{]} Monetary Authority of Singapore. 2026. Safeguards for Agentic
Finance at Runtime (SAFR). White paper, version 1.1, Monetary Authority
of Singapore.
\url{https://www.mas.gov.sg/-/media/mas-media-library/development/fintech/ai-safr/safr-v1-1.pdf}.

\leavevmode\vadjust pre{\hypertarget{ref-CERTIn2026AIAudit}{}}%
{[}41{]} Indian Computer Emergency Response Team. 2026. AI Security
Auditing Guidelines. Guidelines, version 1.0, CERT-In.
\url{https://www.cert-in.org.in/PDF/AI_Security_Auditing_Guidelines.pdf}.

\leavevmode\vadjust pre{\hypertarget{ref-Tang2026SAGEFin}{}}%
{[}42{]} R. Tang, Q. Liu, Y. Zhang, Y. Yang, X. Chen, and C. Dong. 2026.
Context Is Not Authority: Structured Runtime Governance for Financial
Market Agents. arXiv preprint 2608.09025v2.
\url{https://arxiv.org/abs/2608.09025v2}.

\leavevmode\vadjust pre{\hypertarget{ref-Szpruch2026RuntimeGovernance}{}}%
{[}43{]} L. Szpruch, A. Sudjianto, T. Bhatti, and G. Ang. 2026. Scalable
Runtime Governance for Agentic AI in Financial Services. Working paper,
SSRN. \url{https://papers.ssrn.com/abstract=6567199}.

\leavevmode\vadjust pre{\hypertarget{ref-Tan2026SRFraud}{}}%
{[}44{]} X. Tan, Y. Ma, and X. Zhang. 2026. SR-Fraud: An
Outcome-Supervised Reflective LLM Agent Framework for Non-Stationary
Payment Fraud Detection. arXiv preprint 2609.27287v1.
\url{https://arxiv.org/abs/2609.27287v1}.

\leavevmode\vadjust pre{\hypertarget{ref-Sharma2026AuditableFraud}{}}%
{[}45{]} R. Sharma. 2026. Toward Auditable Fraud Detection: Combining
Graph Features, Model Explanations, and Agentic Case Investigation.
arXiv preprint 2607.19266v1. \url{https://arxiv.org/abs/2607.19266v1}.

\leavevmode\vadjust pre{\hypertarget{ref-Stripe2026}{}}%
{[}46{]} Stripe. 2026. ARG Engineering Manager. Stripe Careers.
\url{https://stripe.com/careers/listing/arg-engineering-manager/8113337}.
Accessed 21 October 2026.

\leavevmode\vadjust pre{\hypertarget{ref-StripeARE2026}{}}%
{[}47{]} Stripe. 2026. Abuse Research Engineer. Stripe Careers.
\url{https://stripe.com/careers/listing/abuse-research-engineer/8172503}.
Accessed 21 October 2026.

\leavevmode\vadjust pre{\hypertarget{ref-Sakellariou2022CTIReferenceModel}{}}%
{[}48{]} G. Sakellariou, P. Fouliras, I. Mavridis, and P. Sarigiannidis.
2022. A Reference Model for Cyber Threat Intelligence (CTI) Systems.
\emph{Electronics}, 11(9), Article 1401.
\url{https://doi.org/10.3390/electronics11091401}.

\leavevmode\vadjust pre{\hypertarget{ref-Stripe2026Disputes}{}}%
{[}49{]} Stripe. 2026. How disputes work. Stripe documentation.
\url{https://docs.stripe.com/disputes/how-disputes-work}. Accessed 25
October 2026.

\leavevmode\vadjust pre{\hypertarget{ref-Zhang2026Refund}{}}%
{[}50{]} S. Zhang, E. He, X. Zhan, S. He, R. Xiao, X. Yi, and H. Li.
2026. Generative AI-Enabled Refund Fraud in Chinese E-Commerce:
Investigation on Merchants and Platform Workers. arXiv preprint
2606.03215v1. \url{https://arxiv.org/abs/2606.03215v1}.

\leavevmode\vadjust pre{\hypertarget{ref-Stripe2026MonitoringPrograms}{}}%
{[}51{]} Stripe. 2026. Dispute and fraud card monitoring programs.
Stripe documentation.
\url{https://docs.stripe.com/disputes/monitoring-programs}. Accessed 25
October 2026.

\leavevmode\vadjust pre{\hypertarget{ref-Stripe2026ACHTerms}{}}%
{[}52{]} Stripe. 2026. ACH Payment Terms. Stripe Legal.
\url{https://stripe.com/legal/ACH}. Accessed 25 October 2026.

\leavevmode\vadjust pre{\hypertarget{ref-Thomas2013}{}}%
{[}53{]} K. Thomas, D. McCoy, C. Grier, A. Kolcz, and V. Paxson. 2013.
Trafficking Fraudulent Accounts: The Role of the Underground Market in
Twitter Spam and Abuse. \emph{Proceedings of the 22nd USENIX Security
Symposium}, pp. 195--210.
\url{https://www.usenix.org/conference/usenixsecurity13/technical-sessions/paper/thomas}.

\leavevmode\vadjust pre{\hypertarget{ref-Chen2025Phishing}{}}%
{[}54{]} Z. Chen, D. Luo, Y. Hu, L. Wu, B. He, and Y. Zhou. 2025.
Dissecting Payload-based Transaction Phishing on Ethereum. \emph{Network
and Distributed System Security (NDSS) Symposium 2025}.
\url{https://doi.org/10.14722/ndss.2025.230311}.

\leavevmode\vadjust pre{\hypertarget{ref-Bara2026EpistemicSybil}{}}%
{[}55{]} M. Bara. 2026. Epistemic Sybil Resistance: Multiplying AI
Agents Without Multiplying Evidence. arXiv preprint 2609.01873v1.
\url{https://arxiv.org/abs/2609.01873v1}.

\leavevmode\vadjust pre{\hypertarget{ref-ParasuramanManzey2010}{}}%
{[}56{]} R. Parasuraman and D. H. Manzey. 2010. Complacency and Bias in
Human Use of Automation: An Attentional Integration. \emph{Human
Factors}, 52(3), pp. 381--410.
\url{https://doi.org/10.1177/0018720810376055}.

\leavevmode\vadjust pre{\hypertarget{ref-Wolfe2005Rare}{}}%
{[}57{]} J. M. Wolfe, T. S. Horowitz, and N. M. Kenner. 2005. Rare Items
Often Missed in Visual Searches. \emph{Nature}, 435, pp. 439--440.
\url{https://doi.org/10.1038/435439a}.

\leavevmode\vadjust pre{\hypertarget{ref-Greshake2023}{}}%
{[}58{]} K. Greshake, S. Abdelnabi, S. Mishra, C. Endres, T. Holz, and
M. Fritz. 2023. Not What You've Signed Up For: Compromising Real-World
LLM-Integrated Applications with Indirect Prompt Injection.
\emph{Proceedings of the 16th ACM Workshop on Artificial Intelligence
and Security}, pp. 79--90. Cited version: arXiv 2302.12173v2 (5 May
2023). \url{https://doi.org/10.1145/3605764.3623985}.

\leavevmode\vadjust pre{\hypertarget{ref-Debenedetti2024}{}}%
{[}59{]} E. Debenedetti, J. Zhang, M. Balunovic, L. Beurer-Kellner, M.
Fischer, and F. Tramèr. 2024. AgentDojo: A Dynamic Environment to
Evaluate Prompt Injection Attacks and Defenses for LLM Agents.
\emph{Advances in Neural Information Processing Systems 37, Datasets and
Benchmarks Track}, pp. 82895--82920.
\url{https://doi.org/10.52202/079017-2636}.

\leavevmode\vadjust pre{\hypertarget{ref-Kumar2026}{}}%
{[}60{]} A. A. Kumar. 2026. Loopjacking: Hijacking Human-in-the-Loop
Approval. arXiv preprint 2609.21081v1.
\url{https://arxiv.org/abs/2609.21081v1}.

\leavevmode\vadjust pre{\hypertarget{ref-InvariantLabs2025}{}}%
{[}61{]} L. Beurer-Kellner and M. Fischer. 2025. MCP Security
Notification: Tool Poisoning Attacks. Invariant Labs.
\url{https://invariantlabs.ai/blog/mcp-security-notification-tool-poisoning-attacks}.
Accessed 19 October 2026.

\leavevmode\vadjust pre{\hypertarget{ref-Dash2026}{}}%
{[}62{]} P. Dash, T. Ge, A. Jain, T. Shah, and Z. Shang. 2026. From
Untrusted Input to Trusted Memory: A Systematic Study of Memory
Poisoning Attacks in LLM Agents. arXiv preprint 2606.04329v2.
\url{https://arxiv.org/abs/2606.04329v2}.

\leavevmode\vadjust pre{\hypertarget{ref-Nasr2025}{}}%
{[}63{]} M. Nasr, N. Carlini, C. Sitawarin, S. V. Schulhoff, J. Hayes,
M. Ilie, J. Pluto, S. Song, H. Chaudhari, I. Shumailov, A. Thakurta, K.
Y. Xiao, A. Terzis, and F. Tramèr. 2025. The Attacker Moves Second:
Stronger Adaptive Attacks Bypass Defenses Against LLM Jailbreaks and
Prompt Injections. arXiv preprint 2510.09023v1.
\url{https://arxiv.org/abs/2510.09023v1}.

\leavevmode\vadjust pre{\hypertarget{ref-Stripe2026ThreeDS}{}}%
{[}64{]} Stripe. 2026. Authenticate with 3D Secure. Stripe
documentation.
\url{https://docs.stripe.com/payments/3d-secure/authentication-flow}.
Accessed 19 October 2026.

\leavevmode\vadjust pre{\hypertarget{ref-Stripe2026SEPADebit}{}}%
{[}65{]} Stripe. 2026. SEPA Direct Debit payments. Stripe documentation.
\url{https://docs.stripe.com/payments/sepa-debit}. Accessed 25 October
2026.

\leavevmode\vadjust pre{\hypertarget{ref-Stripe2026CAE}{}}%
{[}66{]} Stripe. 2026. Customer account evaluations. Stripe
documentation.
\url{https://docs.stripe.com/radar/customer-account-evaluations}.
Accessed 19 October 2026.

\leavevmode\vadjust pre{\hypertarget{ref-FTC2020FirstData}{}}%
{[}67{]} Federal Trade Commission. 2020. FTC v. First Data Merchant
Services LLC, No. 1:20-cv-03867 (S.D.N.Y.), Complaint, ECF 1. US
District Court, Southern District of New York. Filed complaint; related
stipulated order and nonadmission terms identified in S1.
\url{https://www.ftc.gov/system/files/documents/cases/first_data_filed_complaint.pdf}.
Accessed 22 October 2026.

\leavevmode\vadjust pre{\hypertarget{ref-FTC2025Paddle}{}}%
{[}68{]} Federal Trade Commission. 2025. FTC v. Paddle.com Market
Limited, No. 1:25-cv-01886 (D.D.C.), Complaint, ECF 1. US District
Court, District of Columbia. Filed complaint; related stipulated order
and nonadmission terms identified in S1.
\url{https://www.ftc.gov/system/files/ftc_gov/pdf/PaddleComplaintForPermanentInjunction,MonetaryJudgment,andOtherRelief.pdf}.
Accessed 21 October 2026.

\leavevmode\vadjust pre{\hypertarget{ref-FTC2019Allied}{}}%
{[}69{]} Federal Trade Commission. 2019. FTC v. AlliedWallet, Inc., No.
2:19-cv-04355 (C.D. Cal.), Complaint, ECF 24. US District Court, Central
District of California. Filed complaint; related stipulated order and
nonadmission terms identified in S1.
\url{https://www.ftc.gov/system/files/documents/cases/alliedwallet_complaint_7-3-19.pdf}.
Accessed 21 October 2026.

\leavevmode\vadjust pre{\hypertarget{ref-Fatima2024}{}}%
{[}70{]} K. Fatima, M. Schuckers, G. Cruz-Ortiz, D. Hou, S. Purnapatra,
T. Andrews, A. Neupane, B. Marshall, and S. Schuckers. 2024. A
Large-Scale Study of Performance and Equity of Commercial Remote
Identity Verification Technologies Across Demographics. \emph{2024 IEEE
International Joint Conference on Biometrics}, pp. 1--8.
\url{https://doi.org/10.1109/IJCB62174.2024.10744432}.

\leavevmode\vadjust pre{\hypertarget{ref-Doerfler2019}{}}%
{[}71{]} P. Doerfler, M. Marincenko, J. Ranieri, Y. Jiang, A. Moscicki,
D. McCoy, and K. Thomas. 2019. Evaluating Login Challenges as a Defense
Against Account Takeover. \emph{Proceedings of the 2019 World Wide Web
Conference}, pp. 372--382.
\url{https://doi.org/10.1145/3308558.3313481}.

\leavevmode\vadjust pre{\hypertarget{ref-Stripe2026PAYG}{}}%
{[}72{]} Stripe. 2026. Pay-as-you-go abuse evaluation. Stripe
documentation. \url{https://docs.stripe.com/radar/pay-as-you-go-abuse}.
Accessed 19 October 2026.

\leavevmode\vadjust pre{\hypertarget{ref-Meltzer2025}{}}%
{[}73{]} J. Meltzer and V. Chadalapaka. 2025. How Stripe Radar responded
to a new wave of card testing. Stripe blog.
\url{https://stripe.com/blog/how-stripe-radar-responded-to-a-new-wave-of-card-testing}.
Accessed 21 October 2026.

\leavevmode\vadjust pre{\hypertarget{ref-Tian2018}{}}%
{[}74{]} H. Tian, S. M. Gaffigan, D. S. West, and D. McCoy. 2018.
Bullet-Proof Payment Processors. \emph{2018 APWG Symposium on Electronic
Crime Research}, pp. 1--11.
\url{https://doi.org/10.1109/eCrime.2018.8376208}.

\leavevmode\vadjust pre{\hypertarget{ref-Xiong2026MemoryManagement}{}}%
{[}75{]} Z. Xiong, Y. Lin, W. Xie, P. He, Z. Liu, J. Tang, H. Lakkaraju,
and Z. Xiang. 2026. How Memory Management Impacts LLM Agents: An
Empirical Study of Experience-Following Behavior. \emph{Proceedings of
the 64th Annual Meeting of the Association for Computational Linguistics
(Volume 1: Long Papers)}, pp. 623--645.
\url{https://doi.org/10.18653/v1/2026.acl-long.27}.

\leavevmode\vadjust pre{\hypertarget{ref-Visa2025VAMP}{}}%
{[}76{]} Visa. 2025. Visa Acquirer Monitoring Program Overview. Visa.
\url{https://corporate.visa.com/content/dam/VCOM/corporate/visa-perspectives/security-and-trust/documents/visa-acquirer-monitoring-program-fact-sheet-2025.pdf}.
Accessed 25 October 2026.

\leavevmode\vadjust pre{\hypertarget{ref-MITRE2026F3}{}}%
{[}77{]} MITRE Center for Threat-Informed Defense. 2026. Fight Fraud
Framework (F3), version 1.2. GitHub,
center-for-threat-informed-defense/fight-fraud-framework. Cited version:
commit e03c3dc (5 November 2026); version 1.2 released 29 September
2026.
\url{https://github.com/center-for-threat-informed-defense/fight-fraud-framework/tree/e03c3dc0fd765e782f35a6293cc8bb340aa14842}.
Accessed 9 November 2026.

\end{CSLReferences}
\end{document}
